\documentclass[reqno]{amsart}
\usepackage{a4wide}
%usepackage{hyperref}

\AtBeginDocument{{\noindent\small
\emph{Electronic Journal of Differential Equations},
Vol. 2026 (2026), No. 08, pp. 1--18.\newline
ISSN: 1072-6691. URL: https://ejde.math.txstate.edu, DOI: 10.58997/ejde.2026.08}
\thanks{\copyright 2026. This work is licensed under a CC BY 4.0 license.}
\vspace{8mm}}

\begin{document}
\title[\hfilneg EJDE-2026/08\hfil
Ground states for Schr\"odinger systems]
{On nodal ground states for Schr\"odinger systems}

\author[A. L. Amadori \hfil EJDE-2026/08\hfilneg]
{Anna Lisa Amadori}

\address{Anna Lisa Amadori \newline
Dipartimento di Scienze e Tecnologie,
Universit\`a di Napoli ``Parthenope'',
Centro Direzionale di Napoli,
Isola C4, 80143 Napoli, Italy}
\email{annalisa.amadori@uniparthenope.it}

\dedicatory{Communicated by Giovanni Molica Bisci }

\thanks{Submitted December 28, 2025. Published January 28, 2026.}
\subjclass[2020]{35J50, 35B05}
\keywords{Coupled nonlinear Schr\"odinger equation; nodal and seminodal solution;
\hfill\break\indent least energy solution; simmetry breaking}

\begin{abstract}
In this article we characterize the least energy nodal and semi-nodal solutions
to some  Schr\"{o}dinger system  as the  minimum on constrained Nehari sets of
codimension 4 and 3, respectively;
thus allowing to compute their Morse index and the exact number of nodal domains.
Next the focus is on the symmetry properties of the sign-changing solutions.
We show that, even though the domain is a ball, ground states are not radial,
and produce other non-radial solutions with the given symmetry.
\end{abstract}

\maketitle
\numberwithin{equation}{section}
\newtheorem{theorem}{Theorem}[section]
\newtheorem{lemma}[theorem]{Lemma}
\newtheorem{corollary}[theorem]{Corollary}
\newtheorem{proposition}[theorem]{Proposition}
\allowdisplaybreaks

\section{Introduction}

We study the Schr\"odinger system
\begin{equation}\label{S}
 \begin{gathered}
 -\Delta u +\lambda_1 u= \mu_1 |u|^{2q-2} u +\beta |u|^{q-2}|v|^q u \quad
  \text{in } \Omega , \\
 -\Delta v +\lambda_2 v= \mu_2 |v|^{2q-2} v +\beta |u|^{q}|v|^{q-2} v \quad
  \text{in } \Omega ,\\
  u=v=0 \quad \text{on } \partial\Omega,
 \end{gathered}
\end{equation}
where $\lambda_1,\lambda_2\ge 0$, $\mu_1,\mu_2>0$, $\Omega$ is a bounded smooth domain
in $\mathbb{R}^N$ ($N=2,3$), and  $2\le q<q_N$, ($q_2=\infty$, $q_3=3$).
The parameter $\beta$ measures the coupling of the system: we are interested in
$\beta<0$ (competing system) or $\beta >0$ but ``small''
 (weakly cooperative system).

A Schr\"{o}dinger system of type  \eqref{S}  arises as a model for various physical
phenomena, in particular in the study of standing
waves for a mixture of Bose-Einstein condensates of two states, see
for example \cite{BoseEinstein}.
From a mathematical viewpoint, it has been so extensively studied that narrowing down
the most significant works is necessarily reductive.
The most extensively studied case is $q=2$, also known as Gross-Pitaewskij equation,
in particular concerning existence and multiplicity of positive solutions;
we mention \cite{LW_AIHP05,BW_JDE06,MMP_JDE06}, among others.
More recently, sign-changing and semi-nodal solutions have attracted the attention of
many mathematicians. When $\beta >0$ is sufficiently small,  radially symmetric
sign-changing solutions with any given number of nodal domains  were constructed
in \cite{MMP_CCM08} on $\mathbb{R}^N$ for $q=2$, and  a similar result for $N=3$ and $2\le q <3$
was proved in \cite{Z_BVP24}.
As for bounded (non spherically symmetric) domains,  the existence of infinitely many
sign-changing and seminodal solutions has been proved in \cite{CLZ_JDE13}
in the weakly cooperative regime, and in \cite{CLZ_ASNSP16} and \cite{LLW_CV15}
in the competitive regime.
These contributions address the case
$q=2$ and investigate  appropriate constrained problems by means of the notion of
vector genus introduced by \cite{TT_AIHP12}.
We also mention  \cite{SW_DCDS15,CS_CCM23,Song_JDE25} for some related systems
of $d\ge 2$ equations.

Here we focus on least energy nodal and semi-nodal solutions, and we characterize
them as constrained minima on a Nehari type set of codimension 4 and 3, respectively.
The solutions display the usual inf-sup characterization of mountain-pass type solutions,
with the paths interpreted here as 4- (respectively 3-) dimensional surfaces rather
than curves.
To manage the complexity of the four (or three) dimensions, we impose the additional
condition $q\ge2$. In that way, the problem is subcritical only for $N\le 3$.
A similar assumption was also made in \cite{MMP_JDE06, Z_BVP24}.
Even though some technical details differ if $\beta<0$ or $\beta>0$,
the competitive and weakly cooperative regimes can be handled simultaneously.

This approach allows to easily compute the number of nodes and the Morse index
of least energy solutions, in particular we see that nodal least energy solutions have
Morse index 4, and each component has exactly two nodal zones
(Proposition \ref{prop:Morseindex}). Likewise, semi-nodal solutions have Morse index 3
 and their components have respectively two and one nodal zones
 (Proposition \ref{prop:seminodal2}).
Furthermore, the Nehari construction can be adapted to subspaces of functions with
pre-assigned symmetries: we give here two examples.


 When $\Omega$ is a ball, it is easy to show the existence of a pure vector
 (resp., a nodal or a semi-nodal) solution with radial symmetry,  which has the
least energy among all radial and nontrivial (resp., nodal or seminodal) solutions
and has \emph{radial} Morse index 2 (resp., 4 or 3).
Observe that the radial Morse index can be less than the Morse index.
The question is whether this provides new solutions or, rather, a symmetry property
of the least energy solutions.

For the scalar Schrodinger equation corresponding to $\beta=0$, this fact is well
understood. The least energy solution on the ball (which has fixed sign) is radial
by the celebrated result by Gidas, Ni and Niremberg \cite{GNN}.
Conversely, comparing the information on the Morse index one sees that any sign
changing radial solution  is not the least energy one in the whole set of sign
changing solutions, see \cite{AftPac}.
Though, the least energy nodal solution inherits some symmetry property of the
domain, precisely it is foliated Schwarz symmetric, because its Morse index is not
greater than $N$, see \cite{PacWet}.

For cooperative Schr\"odinger  systems, a similar sufficient condition for Schwarz
symmetry has been established in \cite{DamPac}: it applies to pure vector ground states,
but not to sign changing solutions in dimension $2$ and $3$.
The estimation of the Morse index of radial solutions  is, in turn, a delicate issue,
and warrants further study.
We limit ourselves here to a perturbative approach, which provides a result concerning
small interaction parameter.

\begin{theorem}\label{thm:symmetry_breaking}
Let $\Omega$ be a ball in dimension $N=2$ or $3$, and $2\le q<q_N$.
There exist $b_1  , b_2> 0$ such that sign-changing (resp., semi-nodal) ground states
are not radial when $|\beta|<b_1$ (resp., $|\beta|<b_2$).
\end{theorem}

If ${\mathrm g}$ is any subgroup of the orthogonal group $O(N)$ one can address to the
subspace of $\mathrm g$-invariant functions
and easily prove the existence of pure vector, nodal and semi-nodal $g$-invariant
solutions which minimize the energy in the respective families.
Though, as any radial function is $g$-invariant, all these $\mathrm g$-invariant
least energy solutions could coincide with the radial ones.
Comparing Morse indices can yield insights into the multiplicity of solutions,
but this relies on knowing not only the precise Morse index of the radial solution,
but also the symmetry features of the associated eigenfunctions.
Studies in this direction have been carried out in the papers \cite{AmaGla20NARWA}
and \cite{GlaIan23Nodea}, concerning the Schr\"odinger equation.
In particular \cite{AmaGla20NARWA}, relying on a singular eigenvalue problem,
provides a decomposition of the spectrum according to spherical coordinates that
naturally extends to systems without substantial modifications,
see also \cite{Ama24JEPE}.
The paper \cite{GlaIan23Nodea}, by means of a fine asymptotic study of radial
solutions to the Lane-Emden problem in the disc, computes the exact Morse index and
deduce a multiplicity result.
Their result  carries over to systems of type \eqref{S} with small interaction
parameter, thanks to a stability argument.

To state the obtained result, we assume that $\Omega$ is a disc in $\mathbb{R}^2$,
write $r$ and $\theta$ for the usual polar coordinates,
\[
H_{k}:=\big\{(u,v)\in \left(H^1_0(\Omega)\right)^2  : \text{$u$ and $v$ are even and
$\frac{2\pi}{k}$-periodic w.r.t. $\theta$} \big\}
\]
for every $k\in \mathbb Z_+$, and call \emph{$k$-symmetric} any function in $H_k$.

\begin{theorem}\label{thm:ksym}
Let $\Omega\subset \mathbb{R}^2$ be a disc and $\lambda_1=\lambda_2=0$. There exists
$q^*\ge 2$ such that for every $q\ge q^*$ and $|\beta|< \bar\beta(q)$ the sign-changing
$k$-symmetric least energy solution is not radial for $k=1,\dots 5$.
\end{theorem}

We also mention that in higher dimension, and for $\beta<0$, nonradial solutions
with a different kind of symmetry have been produced in \cite{CS_CCM23}.

The paper is organized as follows. In Section \ref{sec1} we describe the variational
structure and recollect some facts about ground state solutions (with fixed sign).
Next in Section \ref{sec2} we introduce the constrained nodal Nehari set
\eqref{def:Neharinod}, describe its geometric properties and prove the existence
(Theorem \ref{thm:existence_ nodal}) and some further properties
(Subsection \ref{sub:properties_lesc}) of sign-changing solutions, including the
computation of their Morse index and of the number of nodes, in full details;
moreover in Subsection \ref{sub:seminodal} we apply similar reasoning to produce a
semi-nodal least energy solution.
Section \ref{sec3} is devoted to symmetric solutions and proves Theorems
\ref{thm:symmetry_breaking} and \ref{thm:ksym}.


\section{Preliminary remarks on the variational setting and fixed-sign
solutions}\label{sec1}

Henceforth we let $\Omega$ be a bounded smooth domain in $\mathbb{R}^N$, $N\ge 2$ and
\[
q_N= \begin{cases} +\infty & \text{for } N=2,\\
\frac{N}{N-2} &\text{for } N\ge 3 .
\end{cases}
\]
For $i=1,2$ let $E_1=H^1_0(\Omega)$ equipped with the inner product
\[
\langle w,\phi \rangle_{E_i}
= \int_{\Omega} \nabla w \nabla \phi  +\lambda_i \int_{\Omega} w \phi
 \]
and the induced norm $ \|w\|_{E_i}^2=\langle w,w\rangle_{E_i}$,
which is equivalent to the standard one.
Both $E_1$ and $E_2$ compactly embed into $L^{2q}$ for $1<q<q_N$.
We denote by  $E$ the product space $E_1\times E_2$ with the natural inner product
and norm.

Problem \eqref{S} has a variational structure and the related energy is
$\mathcal I: E \to \mathbb{R}$,
\begin{equation} \label{energy}
\mathcal{I}(u,v) = \int_{\Omega} |\nabla u|^2 + | \nabla v|^2 + \lambda_1 u^2
+ \lambda_2 v^2   - \frac{1}{2q}\int_{\Omega} \mu_1 |u|^{2q}
+ \mu_2 |v|^{2q} + 2\beta |uv|^q .
\end{equation}
Note that weak solutions are critical points for $\mathcal I$.
A ground state solution attains the minimum of $\mathcal I$ among all solutions
(except the trivial one $u=v=0$), i.e.\ realizes
\[
c_{\rm full}=\inf\{ \mathcal I(u,v)  :  (u,v)\in E\setminus\{0\}, \,
 \mathcal I'(u,v)=0\}.
 \]

As $1< q< q_N$, it is completely standard producing ground state solutions by
minimizing $\mathcal I$ on the Nehari set
\begin{equation}\label{def:Nehari0}
 {\mathcal N}_{\rm full} = \big\{ (u,v)\in E\setminus\{0\} :
   I'(u,v) \, (u,v) =0 \big\}.
\end{equation}
But this procedure can lead to a so called semi-trivial solution of type
$(u,0)$ or $(0,v)$, where $u$ and $v$ are the least energy  solutions of the
single equations obtained letting $\beta=0$.
It happens, for instance, if the parameter $\beta$ is negative or positive but small,
see for instance \cite{BW_JDE06}.

In this range, though, nontrivial or ``pure vector'' solutions, i.e.\
 with $u\neq 0, v\neq 0$, do exist.
In particular, there is a nontrivial ground state, which attains
\begin{equation}\label{def:c}
  c=\inf\{ \mathcal I(u,v) : (u,v)\in E, \, u, v \neq 0,  \ \mathcal I'(u,v)=0\} ,
\end{equation}
and can be characterized as the minimum in a constrained Nehari set
\begin{gather} \label{def:gamma}
 \gamma= \inf\left\{ \mathcal I(u,v) : (u,v)\in \mathcal N\right\} , \\
\label{def:Neharivec}
\mathcal N=   \{ (u,v)\in E: u,v\neq 0, \ \mathcal I'(u,v)(u,0)= \mathcal I'(u,v)
(0,v)=0 \}.
\end{gather}


\begin{proposition}\label{prop:existence_vec}
Let $N=2,3$ and $2\le q<q_N$. There exists $\bar\beta>0$ such that
$c=\gamma$  is attained by a function $(u,v)\in\mathcal N$ which solves \eqref{S},
for every  $\beta<\bar \beta$.
Moreover
\[
c=\gamma=\inf_{(u,v)\in \mathcal A }  \sup_{s,t\ge 0}
\mathcal I (s^{\frac{1}q}u, t^{\frac{1}q}v),
\]
with
\[
 \mathcal A = \begin{cases}
\big\{ (u,v) \in E  :  \mu_1 \|u\|_{4}^4 \|v\|_{E_2}^2 > \beta \|uv\|_2^2 \|u\|_{E_1}^2 ,\\
 \quad   \mu_2 \|v\|_{4}^4 \|u\|_{E_1}^2 > \beta \|uv\|_2^2 \|v\|_{E_2}^2 \big\}
& \text{if $q=2$ and $0<\beta<\bar\beta$},
\\[3pt]
E\setminus\{0\} & \text{otherwise.}
\end{cases}
 \]
\end{proposition}

In the case $q=2$, the proof of Proposition \ref{prop:existence_vec} consists in
a mere rearrangement of the arguments employed in \cite[Section 2]{LW_AIHP05}.
Modifications are necessary to handle the case $q>2$, that we do not report here
because they are explained in detail in the following section about sign-changing
solutions.

Starting from the  characterization of the ground state as a minimum on the Nehari set,
one can compute the Morse index and  the number of nodal zones.
We recall that, as the energy functional $\mathcal I$ is of class $C^2$ on $E$,
the Morse index of any solution $(u,v)$ is the maximal dimension of a subspace of
$E$ where the quadratic form related to $ \mathcal{I}''(u,v)$
is negative defined.

\begin{proposition}\label{prop:Morse_vec}
For $q\in [2,q_N)$ and $\beta <\bar\beta $, a nontrivial ground state has Morse
index 2.
\end{proposition}

\begin{proof}
The set $\mathcal N$ is a regular manifold of codimension $2$ in $E$, hence the Morse
index of any solution is at most $2$.
On the other hand, one sees that the Morse index is at least $2$ by showing that
the quadratic form associated to $\mathcal I''(u,v)$ is negative on the 2-dimensional
space spanned by $(u,0)$ and $(0,v)$.
 For every $(t,s)\in \mathbb{R}^2$ $(t,s)\neq(0,0)$ we have
\begin{align*}
 \langle \mathcal I''(u,v)  (tu,sv) , (tu, sv)\rangle
&= t^2 \left( \|u\|_{E_1}^2 -(2q\!-\!1)\mu_1 \|u\|_{2q}^{2q}
  - \beta (q\!-\!1)\|uv\|_q^q \right)  \\
&\quad + s^2\left( \|v\|_{E_2}^2 -(2q\!-\!1)\mu_2 \|v\|_{2q}^{2q}
  - \beta (q\!-\!1)\|uv\|_q^q \right)  - 2\beta q ts \|uv\|_q^{q}\\
& = -2(q-1) \left( t^2 \|u\|_{E_1}^2 + s^2 \|v\|_{E_2}^2\right)
 +\beta q (t-s)^2 \|uv\|_q^{q},
\end{align*}
 since $(u,v)$ solves \eqref{S}.

When $\beta\le 0$ the above quantity certainly is negative.
When $\beta>0$, instead, using again that $(u,v)$ solves \eqref{S} we write
\begin{align*}
&\langle \mathcal I''(u,v)  (tu,sv) , (tu, sv)\rangle \\
&= -2(q-1) \int_{\Omega}\left(\mu_1 t^2 |u|^{2q} + \mu_2 s^2|v|^{2q} \right)
 -\left(\beta (q-2) (t^2+s^2) +2\beta qts \right) \int_{\Omega}|uv|^q \\
&=  -2(q-1) \int_{\Omega}\left(\sqrt{\mu_1} |t| |u|^{q} - \sqrt{\mu_2} |s||v|^{q}
 \right)^2 \\
&\quad - \beta (q-2)(t-s)^2\int_{\Omega}|uv|^q
  - 4(q-1)\left(\sqrt{\mu_1\mu_2} |ts|+\beta ts \right) \int_{\Omega}|uv|^q
<0 ,
\end{align*}
because $\beta<\bar\beta\le \sqrt{\mu_1\mu_2}$.
\end{proof}

We remark that for large positive values of $\beta$ the ground state solution
which realizes the minimum on $\mathcal N_{\rm full}$ is nontrivial and has
Morse index 1, see \cite{MMP_JDE06}.

Knowing the Morse index, some qualitative properties follows.
First, the components of a least energy nontrivial solution  have fixed sign
(note that if $(u,v)$ is any solution to \eqref{S} then also $(\pm u, \pm v)$
solve \eqref{S}).

\begin{corollary}\label{cor:positive}
For $q\in [2,q_N)$ and $\beta <\bar\beta $, a  nontrivial ground state $(u,v)$
is (component-wise) sign definite, i.e. $\pm u>0$ and $\pm v>0$ on $\Omega$.
\end{corollary}

\begin{proof}
Assume by contradiction that neither $u^+$ nor $u^-$ is identically zero.
Then the quadratic form associated to $\mathcal I''(u,v)$ is negative defined on
the three-dimensional space spanned by $(u^+, 0)$, $(u^-, 0)$, $(0,v)$.
The computations are very similar to the ones in the proof of Proposition
\ref{prop:Morse_vec} and we do not repeat them (see also Proposition
\ref{prop:Morseindex} later on). But this contradicts the fact that $(u,v)$ has
Morse index 2.
Hence $\pm u, \pm v \ge 0$, and using Hopf boundary Lemmas it is easy to show that
$\pm u, \pm v >0$  on $\Omega$.
\end{proof}

When the set  $\Omega$ is radially symmetric, a ball or an annulus, then the
nontrivial ground state inherits some symmetry property of the domain, at least in
the cooperative regime.
Recall that $(u,v)$ is said \emph{foliated Schwarz symmetric} if there exist a
unitary vector $p$ in $\mathbb{R}^N$ such that $u(x)$ and $v(x)$ depend only on $r = |x|$ and
$\varphi = \arccos (x\cdot p /|x|)$ and are nonincreasing in $\theta$.
Applying \cite[Theorem 1.1]{DamPac} gives the following result.

\begin{corollary}\label{cor:foliated}
Under the assumptions of Proposition \ref{prop:existence_vec}, if $\beta > 0$,
then a  nontrivial ground state is  foliated Schwarz symmetric.
Moreover if it is not radial, it is (component-wise) strictly decreasing in the
angular variable.
\end{corollary}

See also the example in \cite[Section 4]{DamPac}.

\section{Nodal solutions}\label{sec2}

Given a function $w$ we write $w^+$ and $w^-$ for its positive and negative part,
i.e. $w^{\pm}=\max\{\pm w, 0\}$.
We will say that $(u,v)$ is  (component-wise) sign-changing, or nodal,
if $u^{\pm}$ and $v^{\pm}$ are all non-trivial, and
write $E_{\rm nod}$ for the subset of $E$ made up by sign-changing functions, i.e.
\[
E_{\rm nod} = \{(u,v) \in E : u^{\pm},v^{\pm}\neq0\}.
\]
A function $(u,v)\in E_{\rm nod}$ which solves \eqref{S} is said a nodal least energy
solution if $\mathcal I(u,v)$ is equal to
\[
c_{\rm nod}= \inf \left\{ \mathcal I(u,v) : (u,v)\in E_{\rm nod},
 \mathcal I'(u,v) =0\right\}.
\]
One can not expect to produce a nodal neither a semi-nodal  solution by minimization
in the  standard nodal Nehari set
\[
\mathcal M = \{ (u,v)\in E : (u^+,v^+)\neq 0, \, (u^-,v^-)\neq 0, \,
 \mathcal I'(u,v)(u^+,v^+)= \mathcal I'(u,v) (u^- ,v^-)=0\} .
 \]
Indeed, if $(\bar u, \bar v)$ is a nontrivial least energy solution, then
$(|\bar u| , -|\bar v|)$ is itself a least energy solution and belongs  to $\mathcal M$,
but none of  its components are sign-changing.

Here we produce a nodal least energy solution by minimizing the energy functional on
a constrained set of codimension 4.
This section is organized as follows. First we introduce the nodal Nehari set and
describe its geometrical properties, next in Subsection \ref{sub:existence_lesc}
we adapt the standard Nehari technique to produce a least energy nodal solution,
and in Subsection \ref{sub:properties_lesc} we deduce some further properties.
At last we illustrate how to obtain similar results for semi-nodal solutions in
Subsection \ref{sub:seminodal}.

\subsection{Nodal Nehari set}
We define the nodal Nehari set as
\begin{equation} \label{def:Neharinod}
 \begin{aligned}
 \mathcal N_{\rm nod}&:=  \big\{ (u,v)\in E_{\rm nod}  :
 \mathcal I'(u,v)(u^{\pm},0)= \mathcal I'(u,v) (0,v^{\pm})=0 \big\}\\
 &=  \big\{ (u,v)\in E_{\rm nod}  : \|u^{\pm}\|_{E_1}^2= \mu_1 \|u^{\pm}\|_{2q}^{2q}
  + \beta \|u^{\pm} v\|_{q}^q,
  \|v^{\pm}\|_{E_2}^2= \mu_2 \|v^{\pm}\|_{2q}^{2q} + \beta \|u v^{\pm} \|_{q}^q \big\}.
\end{aligned}
\end{equation}
To describe the geometric properties of $\mathcal N_{\rm nod}$,
we fix  $(u,v)\in  E_{\rm nod}$ and define the function $\theta:[0,\infty)^4 \to \mathbb{R}$
\begin{equation} \label{def:theta}
\begin{aligned}
 \theta(a,b,c,d)
& = \mathcal I (a^{\frac{1}{q}} u^+ -b^{\frac{1}q} u^-, c^{\frac{1}q} v^+ -d^{\frac{1}q} v^-)
\\
& = \frac{1}2  W \cdot (a^{\frac{2}{q}} , b^{\frac{2}q} ,
c^{\frac{2}q} ,d^{\frac{2}q} ) - \frac{1}{2q} M (a,b,c,d) \cdot (a,b,c,d)
\end{aligned}
\end{equation}
where
\begin{gather}\label{def:W}
 W=W(u,v)= \left( \|u^+\|_{E_1}^2 , \|u^-\|_{E_1}^2 , \|v^+\|_{E_2}^2 ,
 \|v^-\|_{E_2}^2 \right), \\
\label{def:M}
 M=M(u,v)= \begin{pmatrix}
 \mu_1\|u^+\|_{2q}^{2q} & 0 & \beta \|u^+v^+\|_q^q & \beta \|u^+v^-\|_q^q\\
 0&  \mu_1\|u^-\|_{2q}^{2q} & \beta \|u^-v^+\|_q^q & \beta \|u^-v^-\|_q^q\\
 \beta \|u^+v^+\|_q^q & \beta \|u^-v^+\|_q^q &  \mu_2 \|v^+\|_{2q}^{2q} & 0 \\
 \beta \|u^+v^-\|_q^q& \beta \|u^-v^-\|_q^q& 0& \mu_2\|v^-\|_{2q}^{2q}
 \end{pmatrix}
\end{gather}

Let
\begin{equation}\label{def:E_0}
 E_0 = \{ (u,v)\in E_{\rm nod} : M(u,v) \text{ is positive defined}\}.
\end{equation}
This set is  nonempty because it contains any couple
$(u,v)\in E_{\rm nod}$ such that $u$ and $v$ have disjoint support. Furthermore the
following holds.

\begin{lemma}\label{lem:M>0}
If $|\beta|<\sqrt{\mu_1\mu_2}/2$, then $E_0=E_{\rm nod}$.
If $\beta \le 0$, then $E_0$ contains any $(u,v)\in E_{\rm nod}$ such that
\begin{equation}\label{eq:E0}
 \mu_1 \|u^{\pm}\|_{2q}^{2q} + \beta \|u^{\pm} v\|_{q}^q>0, \quad
 \mu_2 \|v^{\pm}\|_{2q}^{2q} +\beta \|u v^{\pm} \|_{q}^q >0.
\end{equation}
In particular $\mathcal N_{\rm nod} \subset E_0$.
\end{lemma}

\begin{proof}
For every $(u,v)\in E_{\rm nod}$ and $(a,b,c,d)\in\mathbb{R}^4\setminus\{0\}$ we have
\begin{align*}
&M(u,v) (a,b,c,d) \cdot (a,b,c,d)\\
&= \mu_1 \|u^{+}\|_{2q}^{2q} a^2 + \mu_1 \|u^-\|_{2q}^{2q} b^2 +
  \mu_2 \|v^{+}\|_{2q}^{2q} c^2
  \mu_2 \|v^-\|_{2q}^{2q} d^2
  \\
&\quad +2\beta \|u^+ v^+\|_{q}^q a c  +2\beta \|u^+ v^-\|_{q}^q a d  +2\beta \|u^- v^+\|_{q}^q b c  +2\beta \|u^- v^-\|_{q}^q b c
  \intertext{and  Holder's inequality implies}
&\ge \mu_1 \|u^{+}\|_{2q}^{2q} a^2 + \mu_1 \|u^-\|_{2q}^{2q} b^2 +
  \mu_2 \|v^{+}\|_{2q}^{2q} c^2
  \mu_2 \|v^-\|_{2q}^{2q} d^2
  \\
&\quad -2|\beta| \left( \|u^+\|_{2q}^q \|v^+\|_{2q}^q |a c| + \|u^+\|_{2q}^q \| v^-\|_{2q}^q |a d|  + \|u^- \|_{2q}^q \|v^+\|_{2q}^q |b c|  + \|u^- \|_{2q}^q \|v^-\|_{2q}^q |b d| \right)
  \\
&> \frac{1}{2}
  \left(\sqrt{\mu_1} \|u^{+}\|_{2q}^{q} |a|
  -\sqrt{\mu_2} \|v^{+}\|_{2q}^{2q} |c|  \right)^2
  +\frac{1}{2}
  \left(\sqrt{\mu_1} \|u^{+}\|_{2q}^{q} |a|
  -\sqrt{\mu_2} \|v^{-}\|_{2q}^{2q} |d| \right)^2 \\
&\quad +\frac{1}{2}
  \left(\sqrt{\mu_1} \|u^{-}\|_{2q}^{q} |b|
  -\sqrt{\mu_2} \|v^{+}\|_{2q}^{2q} |c|  \right)^2
  +\frac{1}{2}
  \left(\sqrt{\mu_1} \|u^{-}\|_{2q}^{q} |b|
  -\sqrt{\mu_2} \|v^{-}\|_{2q}^{2q} |d| \right)^2 \ge 0,
 \end{align*}
if $2|\beta|<\sqrt{\mu_1\mu_2}$.
Otherwise,  if $\beta < 0$ and $(u,v)$ satisfies  \eqref{eq:E0}, then
\begin{align*}
M(u,v) (a,b,c,d) \cdot (a,b,c,d)
&>-\beta  \Big(\|u^+ v^+\|_{q}^q (a-c)^2 + \|u^+ v^-\|_{q}^q (a-d)^2
 + \|u^- v^+\|_{q}^q (b-c)^2  \\
&\quad +  \|u^- v^-\|_{q}^q (b-d)^2 \Big)
 \ge 0 .
\end{align*}
\end{proof}

\begin{lemma}\label{prop:theta_qnot2}
Let $q\in (2,q_N)$ and $\beta<\sqrt{\mu_1\mu_2}/2$, or $q=2$ and $\beta \le 0$.
Then for every $(u,v)\in  E_0$  the function
$\theta$ has a unique maximum point in $ (0,\infty)^4$ characterized by the condition
$(a^{\frac{1}q} u^+ -b^{\frac{1}q} u^-, c^{\frac{1}q} v^+ -d^{\frac{1}q} v^-)
 \in \mathcal N_{\rm nod}$.
In particular $(u,v)\in \mathcal N_{\rm nod}$ if and only if $\theta$ achieves its
maximum at $(1,1,1,1)$.
\end{lemma}

\begin{proof}
From the representation in \eqref{def:theta}, and the fact that $M$ is positive defined,
it follows that $\theta(a,b,c,d) \to -\infty$ if $a+b+c+d \to + \infty$. Hence $\theta$
has a maximum point $(a,b,c,d) \in [0,\infty)^4$.
Furthermore $\theta(a,a,a,a)>0=\theta(0,0,0,0)$ for sufficiently small $a$, hence the
maximum point does not fall in the origin.

Let us check that it is an interior point.
If, for instance, $a=0$, then
\[
0\ge \lim_{h\to 0^+} \frac{1}{h}\left[\theta(h, b,c,d) -  \theta(0, b,c,d)\right]
= \lim_{h\to 0^+} \frac{1}{2} \|u^+\|_{E_1}^2 h^{\frac{2-q}q}
- \frac{\beta}{q} \Big[ c \int_{\Omega} |u^+v^+|^q + d \int_{\Omega} |u^+v^-|^q \Big] ,
\]
which implies that either $q <2$ or $q=2$ and $\beta >0$.

So under the present assumptions $\theta$ attains its maximum at some
$(a,b,c,d) \in (0,\infty)^4$, and $\nabla \theta (a,b,c,d)=0$, i.e.
\[
 \left( \|u^+\|_{E_1}^2 a^{\frac{2-q}{q}} , \|u^-\|_{E_1}^2 b^{\frac{2-q}{q}},
 \|v^+\|_{E_2}^2 c^{\frac{2-q}{q}}, \|v^-\|_{E_2}^2 d^{\frac{2-q}{q}}\right)
= M (a,b,c,d) ,
\]
which is equivalent to
$(a^{\frac{1}q} u^+ -b^{\frac{1}q} u^-, c^{\frac{1}q} v^+ -d^{\frac{1}q} v^-)
\in \mathcal N_{\rm nod}$ by a trivial computation.

There are no other  critical points   because $\theta$ is a strictly concave function.
Indeed for every $(a,b,c,d)\in (0,\infty)^4$ we have
\[
- D^2\theta(a,b,c,d) =\frac{q-2}{q^2}
\operatorname{diag}\Big( \|u^+\|_{E_1}^2 a^{\frac{2(1-q)}{q}} ,
\|u^-\|_{E_1}^2 b^{\frac{2(1-q)}{q}}, \|v^+\|_{E_2}^2 c^{\frac{2(1-q)}{q}},
\|v^-\|_{E_2}^2 d^{\frac{2(1-q)}{q}}\Big)
+ \frac{1}q M .
\]
\end{proof}

Note that Proposition \ref{prop:theta_qnot2}  ensures that
$\mathcal N_{\rm nod}$ is not empty.
Proposition \ref{prop:theta_qnot2} does not hold for $q=2$ and $\beta > 0$ because
there are functions $(u,v)\in E_{0}$ for which $\theta$ attains its maximum at the
boundary of $[0,\infty)^4$. We therefore argue in the subset
\begin{equation}\label{def:O}
 \mathcal O = \{ (u,v) \in E_0 :  M^{-1} W \in (0,\infty)^4\} .
\end{equation}
where $M$ and $W$ are defined by \eqref{def:W}, \eqref{def:M}.

\begin{lemma}\label{lem:O}
 The set $\mathcal O$ is nonempty and $\mathcal N_{\rm nod}\subset \mathcal O$.
\end{lemma}

\begin{proof}
Take $w_1, w_2, z_1, z_2$ any smooth nonnegative functions with disjoint supports
and define $w=w_1-w_2$, $z=z_1-z_2$.
Now $M(w,z)$ is diagonal and it is clear that $(w,z)\in \mathcal O$.

Next,  for every $(u,v)\in E_{\rm nod}$ we have
\begin{align*}
M (1,1,1,1)
&= \Big(\mu_1\|u^+\|_{4}^{4} + \beta \|u^+v\|_2^2 ,
 \mu_1\|u^-\|_{4}^{4} + \beta \|u^-v\|_2^2 , \\
&\quad \mu_2 \|v^+\|_{4}^{4} +\beta \|uv^+\|_2^2 ,
\mu_2\|v^-\|_{4}^{4} +\beta \|uv^-\|_2^2 \Big) .
\end{align*}
So $(u,v)\in \mathcal N_{\rm nod}$ if and only if  $M(1,1,1,1)=W$, which implies
$(u,v)\in \mathcal O$.
\end{proof}

\begin{lemma}\label{prop:theta_q2}
Let $q=2$ and $0<\beta<\sqrt{\mu_1\mu_2}/2$.
Then for every $(u,v)\in  \mathcal O$  the function
$\theta$ has a unique maximum point in $ (0,\infty)^4$ characterized by the
condition $(\sqrt{a} u^+ -\sqrt{b} u^-, \sqrt{c} v^+ -\sqrt{d} v^-)
\in \mathcal N_{\rm nod}$.
\end{lemma}

\begin{proof}
Now $\nabla\theta(a,b,c,d)= W- M(a,b,c,d)$, hence for $(u,v)\in \mathcal O$ the
function $\theta $ has an unique critical point  $(a,b,c,d)= M^{-1} W \in (0,\infty)^4$.
Such critical point is a global maximum by concavity.
Furthermore $W-M(a,b,c,d)=0$ is equivalent to
$(\sqrt{a} u^+ -\sqrt{b} u^-, \sqrt{c} v^+ -\sqrt{d} v^-) \in \mathcal N_{\rm nod}$.
\end{proof}

In particular, $\mathcal N_{\rm nod}$ is nonempty also when $q=2$ and
$0<\beta<\sqrt{\mu_1\mu_2}/2$.

\subsection{Existence of a least energy nodal solution}\label{sub:existence_lesc}
Let
\begin{gather}
c_{\rm nod} =\inf \{ \mathcal I(u,v) : (u,v)\in E_{\rm nod} , \  \mathcal I'(u,v)=0\}, \\
\gamma_{\rm nod} =\inf \{ \mathcal I(u,v) : (u,v)\in \mathcal N_{\rm nod} \}.
\end{gather}
Of course $\mathcal N_{\rm nod}$ contains every sign changing solutions of \eqref{S},
hence $\gamma_{\rm nod} \le c_{\rm nod}$. Here we show that
$\gamma_{\rm nod}=c_{\rm nod}$ is attained by a least energy sign changing solution.

\begin{theorem}\label{thm:existence_ nodal}
For every $q\in [2,q_N)$, there exists $\beta_0>0$ such that  for every
$\beta < \beta_0$ the infimum $\gamma_{\rm nod}$
is attained by a (componentwise) sign changing function which solves \eqref{S}.
In particular $c_{\rm nod}=\gamma_{\rm nod}$.
\end{theorem}

By going through the proof, it becomes clear that $\beta_0$  can be taken as
$\frac{1}2\sqrt{\mu_1\mu_2}$ for $q>2$.

\begin{proof}
We begin by noticing that for $(u,v)\in \mathcal N_{\rm nod}$ we have
 \begin{equation} \label{I_Nehari_nod}
 \mathcal I(u,v)=\frac{q-1}{2q}\left[ \mu_1 \|u\|_{2q}^{2q}
 + \mu_2\|v\|_{2q}^{2q}  +2\beta \|uv\|_q^q \right]
 = \frac{q-1}{2q}\left[  \|u\|_{E_1}^{2} + \|v\|_{E_2}^{2} \right]  >0.
 \end{equation}
hence $\gamma_{\rm nod}\ge 0$.
\smallskip

\noindent\textbf{Step 1:} Convergence of a minimizing sequence.
Let $(u_n,v_n)$ a minimizing sequence. By standard compactness arguments, up to a
subsequence,  $(u_n,v_n)$ converges strongly in $(L^{2q}(\Omega))^2$ and weakly in $E$
to some  $(u,v) \in E$.
 In particular $u^{\pm}_n, v^{\pm}_n$ converge strongly in $L^{2q}(\Omega) $ to
 $u^{\pm}, v^{\pm}$ and
\begin{gather}
  \label{conv_I_nod}
  \frac{q-1}{2q}[ \mu_1 \|u\|_{2q}^{2q} + \mu_2\|v\|_{2q}^{2q}
  +2\beta \|uv\|_q^q]
= \lim_{n\to\infty} \mathcal I(u_n,v_n) =\gamma_{\rm nod} ,\\
\label{Fatou1}
  \|u^{\pm}\|_{E_1}^2 \le \lim_{n\to\infty} \|u^{\pm}_n\|_{E_1}^2
  =    \mu_1 \|u^{\pm}\|_{2q}^{2q} + \beta \|u^{\pm}v\|_q^q , \\
\label{Fatou2}
  \|v^{\pm}\|_{E_2}^2 \le \lim_{n\to\infty} \|v^{\pm}_n\|_{E_2}^2 =
 \mu_2 \|v^{\pm}\|_{2q}^{2q} + \beta \|uv^{\pm}\|_q^q .
 \end{gather}
Furthermore there exists a constant $B$ not depending by $\beta$ such that
\begin{equation}\label{bound_nod}
  \|u^{\pm}_n\|_{2q}, \|v^{\pm}_n\|_{2q}\le B
\end{equation}
for every $n$.
To check \eqref{bound_nod}, take $w_1, w_2, z_1, z_2$ any smooth nonnegative functions
with disjoint supports.
It is clear that  $(w_1-w_2, z_1-z_2)$ is in $E_0$, and also in $\mathcal O$ when
$q=2$ and $\beta>0$. Then, let  $w=aw_1-bw_2$, $z=cz_1-dz_2$, where $a,b,c,d$ are
chosen according to Propositions \ref{prop:theta_qnot2}, \ref{prop:theta_q2} in such
a way that $(w,z)\in\mathcal N_{\rm nod}$. So
 \[
\gamma_{\rm nod} \le \mathcal I(w,z)
= \frac{1}2\|(w,z) \|_E^2 - \frac{\mu_1}{2q}\|w\|_{2q}^{2q} -
 \frac{\mu_2}{2q}\|z\|_{2q}^{2q} =B.
 \]
If $\beta \ge 0$, \eqref{bound_nod} readily  follows by  \eqref{conv_I_nod}.
If $\beta \le 0$, instead, using also \eqref{I_Nehari_nod} gives a bound for
$\|u^{\pm}\|_{E_1}$, $\|v^{\pm}\|_{E_2}$ and  \eqref{bound_nod} follows by
Sobolev immersion.
\smallskip

\noindent\textbf{Step 2:} The limit function belongs to $E_0$ for every
$\beta <\frac{1}{2}\sqrt{\mu_1\mu_2}$ if $q>2$.
If $q=2$, there exists $\beta_0\in (0,\frac{1}{2}\sqrt{\mu_1\mu_2}]$ such that the
same holds true  for every $\beta <\beta_0$.

First we  show that $(u,v)\in E_{\rm nod}$: we only check that
$u^+\neq 0$, $u^-,v^{\pm}\neq 0$ follow similarly.
Let $C_1$ be the best constants of the Sobolev embedding of $E_1$ into
$L^{2q}(\Omega)$.
 Now
\begin{align*}
 \frac{1}{C_1} \|u^{\pm}_n\|_{2q}^2 \le \|u^{\pm}_n\|_{E_1}^2
&=\mu_1 \|u^{\pm}_n\|_{2q}^{2q} + \beta \|u^{\pm}_nv_n\|_q^2q \\
&\leq  \mu_1 \|u^{\pm}_n\|_{2q}^{2q} + \beta^+ \|u^{\pm}_n\|_{2q}^q\|v_n\|_{2q}^q
 \quad\text{(by H\"older's inequality)} \\
&\leq \|u^{\pm}_n\|_{2q}^q (\mu_1 \|u_n\|_{2q}^q + 2\beta^+ B ) \quad
 \text{(by \eqref{bound_nod})}.
\end{align*}
where $\beta^+=\max(\beta,0)$.
If $q>2$, then clearly $\|u^{\pm}_n\|_{2q}$ can not vanish.
 But also for $q=2$,  there exists $\beta_0>0$ such that
$\|u^{\pm}_n\|_{4}$ is bounded away from zero for $\beta <\beta_0$.

If $2|\beta|<\sqrt{\mu_2\mu_2}$ there is nothing left to prove.
 Otherwise  \eqref{Fatou1} and \eqref{Fatou2} ensure that $(u,v)\in E_0$,
thanks to Lemma \ref{lem:M>0}.
\smallskip

\noindent\textbf{Step 3:} the limit function  belongs to
$\mathcal N_{\rm nod}$ and achieves its minimum.
The proof differs if $q=2$ or $q>2$, so we split it in two.
\smallskip

\noindent\textbf{Step 4:} For $q=2$,
if $\beta$ is positive it is needed to first check that $(u,v)\in \mathcal O$.
Assume by contradiction that $M^{-1}W \notin (0,\infty)^4$, to fix idea that its
first component is nonpositive. By computations
\begin{equation}\label{noia1}
 M_1 \|u^+\|_{E_1}^2 + M_2 \|u^-\|_{E_1}^2 + M_3 \|v^+\|_{E_2}^2
 + M_4  \|v^-\|_{E_2}^2 \le 0,
\end{equation}
where
\begin{gather*}
\begin{aligned}
M_1&=  \mu_1\mu_2^2 \|u^-\|_4^4\|v^+\|_4^4\|v^-\|_4^4
- \mu_2\beta^2\|v^+\|_4^4\|u^-v^-\|_2^4
- \mu_2\beta^2\|v^-\|_4^4\|u^-v^+\|_2^4  \\
&\ge  \mu_1\mu_2^2 \Big(1-\frac{2\beta^2}{\mu_1\mu_2} \Big)
 \|u^-\|_4^4\|v^+\|_4^4\|v^-\|_4^4 ,
\end{aligned} \\
M_2=  \mu_2\beta^2 \|v^+\|_4^4\|u^+v^-\|_2^2 \|u^-v^-\|_2^2
 + \mu_2\beta^2 \|v^-\|_4^4\|u^+v^+\|_2^2 \|u^-v^+\|_2^2 \ge 0 , \\
\begin{aligned}
M_3 &=  \beta^3 \|u^+v^+\|_2^2 \|u^-v^-\|_2^4 - \beta^3 \|u^+v^-\|_2^2
 \|u^-v^+\|_2^2\|u^-v^-\|_2^2 - \mu_1\mu_2\beta\|u^-\|_4^4\|v^-\|_4^4\|u^+v^+\|_2^2
\\
&\ge  - \mu_1\mu_2 {\beta}\Big(1+\frac{\beta^2}{\mu_1\mu_2}\Big)
 \|u^+\|_4^2\|u^-\|_4^4\|v^+\|_4^2\|v^-\|_4^4,
\end{aligned}\\
\begin{aligned}
 M_4 &=  \beta^3 \|u^+v^-\|_2^2 \|u^-v^+\|_2^4 - \beta^3
  \|u^+v^+\|_2^2 \|u^-v^+\|_2^2\|u^-v^-\|_2^2 - \mu_1\mu_2\beta\|u^-\|_4^4\|v^+
  \|_4^4\|u^+v^+\|_2^2
 \\
&\ge  - \mu_1\mu_2 {\beta}\Big(1+\frac{\beta^2}{\mu_1\mu_2}\Big)
 \|u^+\|_4^2\|u^-\|_4^4\|v^+\|_4^4\|v^-\|_4^2.
\end{aligned}
\end{gather*}
Here we have used repeatedly H\"older inequality. Now \eqref{noia1} implies that
\begin{equation}\label{noia2}
\begin{split}
&\mu_2 \Big(1-\frac{2\beta^2}{\mu_1\mu_2} \Big)
 \|u^-\|_4^4\|v^+\|_4^4\|v^-\|_4^4 \|u^+\|_{E_1}^2 \\
&\le   {\beta}\Big(1+\frac{\beta^2}{\mu_1\mu_2}\Big)
 \|u^+\|_4^2\|u^-\|_4^4\|v^+\|_4^2\|v^-\|_4^2
\left( \|v^-\|_4^2\|v^+\|_{E_2}^2 + \|v^+\|_4^2\|v^-\|_{E_2}^2 \right) .
\end{split}
\end{equation}
Concerning the right-hand side of \eqref{noia2}, using \eqref{Fatou2} and Holder
inequality gives
\[
 \|v^{\pm}\|_{E_2}^2
 \le \mu_2 \|v^{\pm}\|_4^2 \Big(\|v^{\pm}\|_{4}^2
+ \frac{\beta}{\mu_2} \| u\|_4^2 \Big)
 \le \mu_2 \Big(1+ \frac{2\beta}{\mu_2} \Big) B
 \|v^{\pm}\|_4^2 ,
\]
thanks to \eqref{bound_nod}.
Eventually estimating from below the left-hand side of \eqref{noia2} by means of
the Sobolev inequality and simplifying all the repeated terms yields
\begin{equation}\label{noia3}
\frac{1}{C_1} \Big(1-\frac{2\beta^2}{\mu_1\mu_2} \Big)
\le 2{\beta}\Big(1+\frac{\beta^2}{\mu_1\mu_2}\Big)\Big(1+ \frac{2\beta}{\mu_2} \Big) B.
\end{equation}
It leads to a contradiction, possibly after choosing a smaller $\beta_0$.

Afterwards  Proposition \ref{prop:theta_q2}  ensures that there exist $a,b,c,d>0$
such that $(\sqrt{a} u^+ - \sqrt{b}u^-, \sqrt{c} v^+ - \sqrt{d}v^-)
\in \mathcal N_{\rm nod}$, and as $\nabla \theta (a,b,c,d)=0$ we have $M(a,b,c,d)=W$,
that is,
\begin{equation}\label{atante}
\begin{gathered}
\begin{aligned}
& \mu_1 a\|u^+\|_4^4 + \beta c \|u^+v^+\|_2^2 + \beta d \|u^+v^-\|_2^2 \\
&= \|u^+\|_{E_1}^2 \\
&\le \mu_1 \|u^+\|_4^4 + \beta \|u^+v^+\|_2^2  + \beta \|u^+v^-\|_2^2
  \quad \text{(by  \eqref{Fatou1})} ,
\end{aligned}\\
\begin{aligned}
&\mu_1 b\|u^-\|_4^4 + \beta c \|u^-v^+\|_2^2 + \beta d \|u^-v^-\|_2^2  \\
&= \|u^-\|_{E_1}^2 \\
&\leq  \mu_1 \|u^-\|_4^4 + \beta \|u^-v^+\|_2^2   + \beta \|u^-v^-\|_2^2 \quad
\text{(by \eqref{Fatou1})},
\end{aligned}\\
\begin{aligned}
&\mu_2 c\|v^+\|_4^4 + \beta a \|u^+v^+\|_2^2 + \beta b \|u^-v^+\|_2^2 \\
&= \|v^+\|_{E_2}^2  \\
&\leq \mu_2 \|v^+\|_4^4 + \beta \|u^+v^+\|_2^2  + \beta \|u^-v^+\|_2^2 \quad
\text{(by \eqref{Fatou2})},
\end{aligned} \\
\begin{aligned}
&\mu_2 d\|v^-\|_4^4 + \beta a \|u^-v^-\|_2^2 + \beta d \|u^-v^-\|_2^2 \\
& = \|v^-\|_{E_2}^2 \\
&\leq \mu_2 \|v^-\|_4^4 + \beta \|u^+v^-\|_2^2  + \beta \|u^-v^-\|_2^2 \quad
\text{(by \eqref{Fatou2})}.
\end{aligned}
\end{gathered}
\end{equation}
Then
 \begin{align*}
\gamma_{\rm nod}
&\leq  \mathcal I (\sqrt{a} u^+ \!-\!\sqrt{b}u^-\!, \sqrt{c} v^+ \!-\!\sqrt{d} v^-)\\
&=\frac{1}4\Big[\mu_1 a^2\|u^+\|_4^4 +\mu_1 b^2\|u^-\|_4^4
 + \mu_2 c^2\|v^+\|_4^4 + \mu_2 d^2\|v^-\|_4^4 \\
& \quad   +2\beta ac \|u^+ v^+\|_2^2 + 2\beta ad \|u^+ v^-\|_2^2
  + 2\beta bc\|u^- v^+\|_2^2 + 2\beta bd \|u^- v^-\|_2^2 \Big] \\
&\leq \frac{1}{4}  \Big[ a \mu_1 \|u^+\|_4^4  + b \mu_1 \|u^-\|_4^4
 + c \mu_2 \|v^+\|_4^4   + d \mu_2 \|v^-\|_4^4
 + \beta (a+c) \|u^+v^+\|_2^2  \\
&\quad + \beta (a+d) \|u^+v^-\|_2^2
 + \beta (b+c)\|u^-v^+\|_2^2  + \beta (b+d) \|u^-v^-\|_2^2
  \Big] \quad  \text{(by \eqref{atante})} \\
&\leq \frac{1}{4}  \big(\mu_1  \|u\|_4^4 +  2\beta\|uv\|_2^2
 + \mu_2 \|v\|_4^4 \big) \quad  \text{(by \eqref{atante})} \\
&= \gamma_{\rm nod}\quad \text{(by \eqref{conv_I_nod})} .
\end{align*}
It follows that \eqref{atante} are indeed equalities, that is
$M(a,b,c,d)=W= M(1,1,1,1)$
and since the matrix $M$ is nondegenerate, it follows that  $a=b=c=d=1$.
It is thus proved that $(u,v)\in \mathcal N_{\rm nod}$ and
$\mathcal I(u,v)=\gamma_{\rm nod}$.
\smallskip

\noindent\textbf{Step 3:} For $q>2$, let
\begin{equation} \label{def:theta_n_nod}
\begin{split}
  \theta_n(a,b,c,d)
&=  \mathcal I (a^{\frac{1}{q}}u_n^+-b^{\frac{1}{q}}u_n^-,c^{\frac{1}q}v_n^+
-d^{\frac{1}{q}}v_n^-)  \\
&= \frac{1}2 \Big( a^{\frac{2}q}\|u^+_n\|_{E_1}^2 + b^{\frac{2}q}\|u^-_n\|_{E_1}^2
  + c^{\frac{2}q}\|v_n^+\|_{E_2}^2 + d^{\frac{2}q}\|v_n^-\|_{E_2}^2 \Big)
  \\
&\quad - \frac{1}{2q} \left( \mu_1 a^2 \|u_n^+ \|_{2q}^{2q} +\mu_1 b^2 \|u_n^- \|_{2q}^{2q}
 + \mu_2 c^2\|v_n^+\|_{2q}^{2q} +\mu_2 d^2\|v_n^-\|_{2q}^{2q} \right)
  \\
&\quad -\frac{\beta}{q} \Big( ac \|u_n^+v_n^+\|_q^q +  ad \|u_n^+v_n^-\|_q^q
+ bc \|u_n^-v_n^+\|_q^q +  bd \|u_n^-v_n^-\|_q^q \Big) .
\end{split}
\end{equation}
We remark that
 \begin{equation} \label{def:theta_n2_nod}
\begin{split}
&\theta_n(a,b,c,d) \\
&= \frac{1}{2q} \Big[ \mu_1 (q  a^{\frac{2}q} \!-a^2) \|u_n^+\|_{2q}^{2q}
  + \mu_1 (q  b^{\frac{2}q} \!-b^2) \|u_n^-\|_{2q}^{2q}
  +\mu_2 (q  c^{\frac{2}q} \!-c^2) \|v_n^+\|_{2q}^{2q}  \\
&\quad + \mu_2 (q  d^{\frac{2}q} \!-d^2) \|v_n^-\|_{2q}^{2q}
  + ( q a^{\frac{2}q} +  qc^{\frac{2}q} - 2\beta ac)
   \|u_n^+v_n^+\|_q^q  + ( q a^{\frac{2}q} +  qd^{\frac{2}q} - 2\beta ad)
    \|u_n^+v_n^-\|_q^q \\
&\quad  + ( q b^{\frac{2}q} +  qc^{\frac{2}q} - 2\beta bc)  \|u_n^-v_n^+\|_q^q
  + ( q b^{\frac{2}q} +  qd^{\frac{2}q} - 2\beta bd)  \|u_n^-v_n^-\|_q^q \Big]
\end{split}
\end{equation}
because $(u_n,v_n)\in \mathcal N_{\rm nod}$.
From the representation \eqref{def:theta_n2_nod} and the strong convergence of
$u_n, v_n$ in $L^{2q}(\Omega)$, $\theta_n$ converges pointwise to the function
$\bar\theta$ defined by substituting $u^{\pm}$ and $v^{\pm}$ to $u^{\pm}_n$ and
$v^{\pm}_n$ in the law \eqref{def:theta_n2_nod}.
Notice that
\begin{gather*}
\begin{aligned}
\nabla\bar\theta
&= \frac{1}{q} \Big(
 a^{\frac{2-q}q}(\mu_1\|u^+\|_{2q}^{2q} +\beta\|u^+v\|_q^q)
 - ( a\mu_1\|u^+\|_{2q}^{2q} +c \beta \|u^+v^+\|_q^q +d\beta \|u^+v^-\|_q^q), \\
& b^{\frac{2-q}q}(\mu_1\|u^-\|_{2q}^{2q} +\beta\|u^-v\|_q^q)
  - (b \mu_1\|u^-\|_{2q}^{2q} +c\beta \|u^-v^+\|_q^q +d \beta \|u^-v^-\|_q^q),\\
&  c^{\frac{2-q}q}(\mu_2\|v^+\|_{2q}^{2q} +\beta\|uv^+\|_q^q)
 - (c \mu_2\|v^+\|_{2q}^{2q} +a\beta \|u^+v^+\|_q^q +b \beta \|u^-v^+\|_q^q),\\
& d^{\frac{2-q}q}(\mu_1\|v^-\|_{2q}^{2q} +\beta\|uv^-\|_q^q)
 - (d \mu_2\|v^-\|_{2q}^{2q} +a\beta \|u^+v^-\|_q^q +b \beta \|u^-v^-\|_q^q)
 \Big),
\end{aligned} \\
 D^2\bar\theta
 = -\frac{q-2}{q^2} \operatorname{diag}
\begin{pmatrix}
 a^{\frac{2(1-q)}q}(\mu_1\|u^+\|_{2q}^{2q} +\beta\|u^+v\|_q^q) \\
 b^{\frac{2(1-q)}q}(\mu_1\|u^-\|_{2q}^{2q} +\beta\|u^-v\|_q^q) \\
 c^{\frac{2(1-q)}q}(\mu_2\|v^+\|_{2q}^{2q} +\beta\|uv^+\|_q^q) \\
 d^{\frac{2(1-q)}q}(\mu_2\|v^-\|_{2q}^{2q} +\beta\|uv^-\|_q^q)
 \end{pmatrix}
- \frac{1}q M
\end{gather*}
Then  $\bar\theta$ is strictly concave and has an unique maximum point at
$(1,1,1,1)$, with $\bar\theta(1,1,1,1)= \gamma_{\rm nod}$.

Furthermore from the representation \eqref{def:theta_n_nod},  the weak convergence
of $u^{\pm}_n, v^{\pm}_n$ in $L^{2q}(\Omega)$, \eqref{Fatou1} and \eqref{Fatou2},
it follows that
$\bar \theta (a,b,c,d)\ge \theta (a,b,c,d)
=\mathcal I(a^{\frac{1}q}u^+\!-b^{\frac{1}q}u^-,c^{\frac{1}q}v^+\! - d^{\frac{1}q}v^-)$.
If $(a_o,b_o,c_o,d_o)$ is the maximum point of $\theta$, then
\[
\gamma_{\rm nod}\le \mathcal I (a_o^{\frac{1}q}u^-b_o^{\frac{1}q}u^-\!,c_o^{\frac{1}q}v^+-
d_o^{\frac{1}q}v^-)
=\theta(a_o,b_o,c_o,d_o) \le \bar\theta(a_o,b_o,c_o,d_o)
\le \bar \theta(1,1,1,1)=\gamma_{\rm nod}
\]
by \eqref{conv_I_nod}.
Therefore $a_o=b_o=c_o=d_o=1$, which yields at once that $(u,v)\in \mathcal N_{\rm nod}$
and $\mathcal I(u,v)=\gamma_{\rm nod}$.
\smallskip

\noindent\textbf{Step 4:} The limit function solves \eqref{S}.
It is needed to check that $\mathcal I'(u,v)=0$, and
one can not use Lagrange multipliers because $\mathcal N_{\rm nod}$ is not a
differentiable manifold, hence we rely on a deformation argument, see for instance
\cite{CTV_AIHP}.
Assume by contradiction that $\mathcal I'(u,v)\neq 0$, so there is
$\Phi=(\phi, \psi)\in E$ such that $\mathcal I'(u,v) \Phi = -2$.
For $\varepsilon>0$, let
\[
A_{\pm}=\left(1\pm\varepsilon\|u^+\|_{2q}^{-q}\right)^{\frac{q}{2(q-1)}}, \quad
B_{\pm}=\left(1\pm\varepsilon\|u^-\|_{2q}^{-q}\right)^{\frac{q}{2(q-1)}},
\]
$C_{\pm}$ and $D_{\pm}$ defined accordingly with $u$ replaced by $v$,
\[
Q = [A_-,A_+]\times[B_-,B_+]\times[C_-,C_+]\times[D_-,D_+],
\]
$W=(U,V):Q\to E$ defined by the law
\[
W=(U,V)=\Big( a^{\frac{1}q}u^+ - b^{\frac{1}{q}}u^-, c^{\frac{1}q}v^+
- d^{\frac{1}q}v^-\Big) .
\]
By continuity, we can take $\varepsilon$ small so that
\begin{equation}\label{I<-1}
 \mathcal I'(W(a,b,c,d)+r\Phi) \Phi \le -1 \quad
  \text{for all } (a,b,c,d)\in Q \text{ and } r\in [0,\varepsilon].
\end{equation}
Next, we take $\eta:Q\to \mathbb{R}$ a smooth function such that $0\le \eta\le \varepsilon$,
$\eta=0$ on $\partial Q$ and $\eta(1,1,1,1)=\varepsilon$
and define $H : Q \to \mathbb{R}^4$,
\begin{align*}
H&=\Big( \mathcal I'(W+\eta\Phi) ((U+\eta\phi)^+,0),
 \mathcal I'(W+\eta\Phi) ((U+\eta\phi)^-,0),  \\
&\quad  \mathcal I'(W+\eta\Phi) (0,(V+\eta\psi)^+), \mathcal I'(W+\eta\Phi) (0,(V+\eta\psi)^-)
 \Big) .
\end{align*}
We remark that on $\partial Q$ the function $\eta$ vanishes, so
\[
H=\left(qa \partial_a\theta, -qb\partial_b\theta, qc\partial_c\theta,
-qd\partial_d\theta\right) .
\]
In particular for $a=A_+$ we have
\begin{align*}
&q A_+ \partial_a\theta(A_+,b,c,d) \\
&= A_+^{\frac{2}q}\|u^+\|_{E_1}^{2}-\mu_1A_+^2\|u^+\|_{2q}^{2q} -\beta A_+c\|u^+v^+\|_q^q
-\beta A_+d\|u^+v^-\|_q^q\\
 \intertext{and because $(u,v)\in \mathcal N_{\rm nod}$}
&= A_+^{\frac{2}q}\Big[ \mu_1(1-A_+^{\frac{2(q-1)}q}) \|u^+\|_{2q}^{2q}
 -\beta (1-A_+^{\frac{q-2}q}c) \|u^+v^+\|_q^q -\beta (1-A_+^{\frac{q-2}q}d)\|u^+v^-
  \|_q^q\Big] \\
&=A_+^{\frac{2}q}\Big[ -\mu_1 \varepsilon \|u^+\|_{2q}^{q} +\beta (1-A_+^{\frac{q-2}q}c)
 \|u^+v^+\|_q^q +\beta (1-A_+^{\frac{q-2}q}d)\|u^+v^-\|_q^q\Big]
\end{align*}
But
\[
1-A_+^{\frac{q-2}q}c \le  1- C_-
= 1- \left(1-\varepsilon\|v^+\|_{2q}^{-q}\right)^{\frac{q}{2(q-1)}} \le\varepsilon\|v^+\|_{2q}^{-q}
\]
because $q/2(q-1)\le 1$. Similarly $1-A_+^{\frac{q-2}q}d \le \varepsilon\|v^-\|_{2q}^{-q} $
and using Holder inequality we end up with
\[
 q A_+ \partial_a\theta(A_+,b,c,d)
 \le \varepsilon A_+^{\frac{2}q}\|u^+\|_{2q}^{q} \left[ -\mu_1 +2\beta \right] <0
\]
for $\beta <\bar \beta$. Similarly one sees that
 \[
q B_+ \partial_b\theta(a,B_+,c,d), \; q C_+ \partial_c\theta(a,b, C_+,d),
\; q D_+ \partial_d\theta(a,b,c,D_+)< 0
\]
and
\[   q A_- \partial_a\theta(A_-,b,c,d) ,
 \; q B_-\partial_b\theta(a,B_-,c,d),
  \; q C_- \partial_c\theta(a,b, C_-,d),
  \; q D_- \partial_d\theta(a,b,c,D_-) >0 .
\]
In that way, the classical Miranda Theorem \cite{Miranda} ensures that there is
$(a,b,c,d)\in Q$ such that $H(a,b,c,d) =0$, i.e.\
$W+\eta\Phi\in \mathcal N_{\rm nod}$. Eventually
\[
  \mathcal I(u,v) = \gamma_{\rm nod} \le \mathcal I(W+\eta\Phi)
  = \mathcal I(W)+ \int_0^{\eta}  \mathcal I'(W+r\Phi) \Phi dr \le \mathcal I(W) - \eta
\]
thanks to \eqref{I<-1}.
Here we have omitted to specify that $W$ and $\eta$ are computed at $(a,b,c,d)$
to simplify notations.
 Besides $\mathcal I(W(a,b,c,d)) =\theta(a,b,c,d) \le \theta(1,1,1,1)
 = \mathcal I(u,v)=\gamma_{\rm nod}$, therefore
$\eta(a,b,c,d)=0$ and $\theta(a,b,c,d) =\mathcal I(u,v)$.
 But then the uniqueness of the maximum point for $\theta$ in
Propositions \ref{prop:theta_qnot2}, \ref{prop:theta_q2} implies that $a=b=c=d=1$
and then, in turn, the contradiction $\eta(a,b,c,d)=\varepsilon$.
\end{proof}

\subsection{Further properties of the least energy solutions}\label{sub:properties_lesc}

If $\mathcal A$ is any subset of $E$, we introduce the minimax value
\[
m(\mathcal A)=  \inf\big\{ \sup_{a,b,c,d\ge 0} \mathcal I (a^{\frac{1}q}u^+
-b^{\frac{1}q}u^-, d^{\frac{1}q}v^+-d^{\frac{1}q}v^-) : (u,v)\in \mathcal A\big\} .
 \]
Taken together, Theorem \ref{thm:existence_ nodal} and Lemmas
\ref{prop:theta_qnot2}, \ref{prop:theta_q2} yield the following  characterization of
the nodal least energy
\[
 c_{\rm nod}=\gamma_{\rm nod}= \begin{cases}
 m({E_{\rm nod}})& \text{if $q\in (2, q_N)$ and  $|\beta|<\sqrt{\mu_1\mu_2}/2 $
  or $q=2$ and $-\sqrt{\mu_1\mu_2}/2 <\beta\le 0$,} \\
 m({E_0}) & \text{if $q\in [2, q_N)$ and $ \beta \le - \sqrt{\mu_1\mu_2}/2$,}  \\
 m({\mathcal O}) & \text{if $q=2$ and  $0<\beta <\beta_0$.}
\end{cases}
\]
Indeed, it is easy to prove the following result.

\begin{corollary}
We have $c_{\rm nod}=\gamma_{\rm nod}= m(E_{\rm nod})$ for every
$\beta< \sqrt{\mu_1\mu_2}/2 $ if $q\in (2, q_N)$, or for every $\beta\le 0$  if $q=2$.
\end{corollary}

\begin{proof}
It is only needed to address negative values of $\beta$, and it is already known that
$c_{\rm nod}=m(E_0) \ge m(E_{\rm nod})$.
Let $(u_n,v_n)$ be a sequence in $E_{\rm nod}$ and define
$\theta_n(a,b,c,d)=\mathcal I (a^{\frac{1}q}u_n^+-b^{\frac{1}q}u_n^-,
d^{\frac{1}q}v_n^+-d^{\frac{1}q}v_n^-)$.
If $ \sup_{[0,\infty)^4} \theta_n \to m(E_{\rm nod})$, in particular
$\theta_n$ must be bounded and recalling the representation \eqref{def:theta}
one see that the matrix $M(u_n,v_n)$ must be positive defined.
Hence  $(u_n, v_n)\in E_0$, so that $m(E_{\rm nod})\ge m(E_0)$ and the proof is
thereby complete.
\end{proof}

Both the Morse index and the number of nodal zones of least energy nodal solutions can
be computed, though $\mathcal N_{\rm nod}$ is not a manifold in $E$, relying on the
arguments in \cite{BW_JDE06}.
We sketch the proof for the sake of completeness.

\begin{proposition}\label{prop:Morseindex}
Any least energy nodal solution has Morse index 4, and both its components have exactly
two nodal zones, meaning that the supports of $u^{\pm}$ and $v^{\pm}$ are connected.
\end{proposition}

\begin{proof}
Let $G:E_{\rm nod}\to\mathbb{R}^4$,
\[
G(u,v)=\begin{pmatrix}
\mathcal I'(u,v) \, (u^+,0) \\ \mathcal I'(u,v) \, (u^-,0) \\
\mathcal I'(u,v) \, (0,v^+) \\ \mathcal I'(u,v) \, (0,v^-)
 \end{pmatrix},
\]
and $H:=E\cap \left(H^2(\Omega)\right)^2$.
Now $\mathcal N_{\rm nod}\cap H = \{(u,v)\in E_{\rm nod}\cap H : G(u,v)=0\}$ and $G$
is a $C^1$ function on $H$ with
\begin{align*}
&G'(u,v)(\phi,\psi) \\
&=\begin{pmatrix}
\int_{\{u>0\}} -\Delta u \phi + \nabla u \nabla \phi
 + 2\lambda_1 u \phi - q \int_{\Omega}
 \left(2\mu_1|u|^{2q-2}+\beta|u|^{q-2}|v|^q\right)u^+ \phi
 + \beta |uv|^{q-2} uu^+v\psi
\\
\int_{\{u<0\}} -\Delta u \phi + \nabla u \nabla \phi + 2\lambda_1 u \phi
 - q \int_{\Omega} \left(2\mu_1|u|^{2q-2}+\beta|u|^{q-2}|v|^q\right)u^- \phi
 + \beta |uv|^{q-2} uu^-v\psi
\\
\int_{\{v>0\}} -\Delta v \psi + \nabla v \nabla \psi + 2\lambda_2 v \psi
 - q \int_{\Omega} \beta |uv|^{q-2} uvv^+
  \phi+\left(2\mu_2|v|^{2q-2}+\beta|u|^{q}|v|^{q-2}\right)v^+ \psi
\\
\int_{\{v<0\}} -\Delta v \psi + \nabla v \nabla \psi
 + 2\lambda_2 v \psi - q \int_{\Omega} \beta |uv|^{q-2} uvv^-
 \phi+\left(2\mu_2|v|^{2q-2}+\beta|u|^{q}|v|^{q-2}\right)v^- \psi
\end{pmatrix},
\end{align*}
see \cite[Lemma 3.1]{BW_JDE06}.

Furthermore for $(u,v)\in \mathcal N_{\rm nod}\cap H$, $G'(u,v)$ is a surjective
operator from $H$ to $\mathbb{R}^4$. Indeed
\begin{align*}
&\left( G'(u,v)(u^+,v), G'(u,v)(u^-,v), G'(u,v)(u,v^+), G'(u,v)(u,v^-) \right) \\
& = -2(q-1)\begin{pmatrix} \|u^+\|_{E_1}^2 & 0 & 0 & 0 \\
 0 & \|u^-\|_{E_1}^2 & 0 & 0 \\
 0 & 0 & \|v^+\|_{E_2}^2 & 0 \\
 0 & 0& 0& \|v^-\|_{E_2}^2 ,
\end{pmatrix}
\end{align*}
and $u^{\pm}, v^{\pm}$ can be approximated by functions of $H^2(\Omega)$.
In this way $\mathcal N_{\rm nod}\cap H$ is a $C^1$ manifold of codimension 4 in
$H$.

If $(u,v)$ is a least energy nodal solution, then $(u,v)\in \mathcal N_{\rm nod}\cap H$
by elliptic regularity, and by minimality the quadratic form associated to
$\mathcal I'' (u,v)$ is nonnegative on $T$, the tangent space of
$\mathcal N_{\rm nod}\cap H$ at $(u,v)$.
Since $T$ has codimension 4 in $H$ and $H$ is dense in $E$, it follows that the
Morse index is at most 4.

On the other hand, one can see that the Morse index is at least 4 %arguing as in
Proposition \ref{prop:MorseAlto}.
by showing that the same quadratic form is negative on the 4-dimensional  space
spanned by $(u^{\pm},0)$ and $(0,v^{\pm})$.
For every $(a,b,c,d)\in \mathbb{R}^4\setminus\{0\}$  we have
\begin{align*}
&\langle\mathcal I''(au^++bu^-,cv^++dv^-) ,(au^++bu^-,cv^++dv^-) \rangle \\
&=  a^2 \left( \|u^+\|_{E_1}^2 -(2q\!-\!1)\mu_1 \|u^+\|_{2q}^{2q}
- \beta (q\!-\!1)\|u^+v\|_q^q \right) \\
&\quad + b^2\left( \|u^-\|_{E_1}^2 -(2q\!-\!1)\mu_1 \|u^-\|_{2q}^{2q}
 - \beta (q\!-\!1)\|u^-v\|_q^q \right)\\
&\quad + c^2\left( \|v^+\|_{E_2}^2 -(2q\!-\!1)\mu_2 \|v^+\|_{2q}^{2q}
  - \beta (q\!-\!1)\|uv^+\|_q^q \right) \\
&\quad  + d^2\left( \|v^-\|_{E_2}^2 -(2q\!-\!1)\mu_2 \|v^-\|_{2q}^{2q}
 - \beta (q\!-\!1)\|uv^-\|_q^q \right)
 \\
&\quad - 2\beta q \left( ac \|u^+v^+\|_q^{q} + ad \|u^+v^-\|_q^{q} +bc \|u^-v^+\|_q^{q}
+ bd \|u^-v^-\|_q^{q} \right)
 \intertext{and since $(u,v)$ solves \eqref{S}}
&= -2(q-1) \left( a^2 \|u^+\|_{E_1}^2 +  b^2 \|u^-\|_{E_1}^2
 + c^2 \|v^+\|_{E_2}^2 + d^2 \|v^-\|_{E_2}^2\right) \\
&\quad + \beta q \left( (a-c)^2 \|u^+v^+\|_q^{q} + (a-d)^2 \|u^+v^-\|_q^{q}
 + (b-c)^2 \|u^-v^+\|_q^{q} +(b-d)^2 \|u^-v^-\|_q^{q} \right) .
\end{align*}
When $\beta\le 0$ this quantity  certainly is negative.

When $\beta\in (0,\sqrt{\mu_1\mu_2}/2)$, instead, using again that $(u,v)$
solves \eqref{S} we write
\begin{align*}
&\langle\mathcal I''(au^++bu^-,cv^++dv^-) ,(au^++bu^-,cv^++dv^-) \rangle  \\
& = -2(q-1) \int_{\Omega}\mu_1 a^2 |u|^{2q} +\mu_1 b^2 |u^-|^{2q}
 + \mu_2 c^2|v^+|^{2q} + \mu_2 d^2|v^-|^{2q} \\
&\quad - (q-2)\beta \int_{\Omega} (a^2+c^2)|u^+ v^+|^q + (a^2+d^2)|u^+
  v^-|^q + (b^2+c^2)|u^- v^+|^q + (b^2+d^2)|u^- v^-|^q \\
&\quad -2\beta q \int_{\Omega} ac|u^+ v^+|^q + ad|u^+ v^-|^q + bc|u^- v^+|^q
 + bd|u^- v^-|^q \\
&\leq   -(q-1) \int_{\Omega}\left(\sqrt{\mu_1} |a| |u^+|^{q}
 - \sqrt{\mu_2}|c||v^+|^q\right)^2 + \left(\sqrt{\mu_1} |a| |u^+|^{q}
 - \sqrt{\mu_2}|d||v^-|^q\right)^2 \\
&\quad  -(q-1) \int_{\Omega}\left(\sqrt{\mu_1} |b| |u^-|^{q}
 - \sqrt{\mu_2}|c||v^+|^q\right)^2 + \left(\sqrt{\mu_1} |b| |u^-|^{q}
 - \sqrt{\mu_2}|d||v^-|^q\right)^2  \\
&\quad - (q-2)\beta \int_{\Omega} (|a|-|c|)^2|u^+ v^+|^q + (|a|-|d|)^2|u^+ v^-|^q
 + (|b|-|c|)^2|u^- v^+|^q + (|b|-|d|)^2|u^- v^-|^q
  \\
&\quad +2\left(2\beta -(q-1)\sqrt{\mu_1\mu_2} \right)\int_{\Omega} |ac||u^
 + v^+|^q + |ad||u^+ v^-|^q + |bc||u^- v^+|^q + |bd||u^- v^-|^q
 <0 .
 \end{align*}

As for the number of nodal regions of $u$ and $v$, assume that the support of $u^+$
splits into two sets $A_1$ and $A_2$ which are the closure of open disjoint sets.
Then letting $u^+_1$ and $u^+_2$ the restrictions of $u$ to the sets $A_1$ and $A_2$,
respectively, and repeating the previous computations one sees that the quadratic
form $\mathcal I''(u,v)$ is negative  on the 5-dimensional  space spanned by $(u^+_1,0)$,
$(u^+_2,0)$, $(u^-,0)$ and $(0,v^{\pm})$, contradicting the fact that the Morse index
is 4.
\end{proof}

\subsection{Semi-nodal solutions} \label{sub:seminodal}

Least energy semi-nodal solutions can be produced in a very similar way.
We seek for a solution $(u,v)$ with $u$ sign-changing and $v$ sign-definite,
which attains the infimum
\[
 c_{\rm sn} =\inf \{ \mathcal I(u,v) : (u,v)\in E,  \, u^{\pm}\neq 0,
  \,| v|>0 , \; \mathcal I'(u,v)=0  \},
 \]
and works on the Nehari type set
\begin{equation} \label{def:Neharisnod}
  \mathcal N_{\rm sn}:=  \big\{ (u,v)\in E\, :  u^{\pm}, v \neq 0, \;
   \mathcal I'(u,v)(u^{\pm},0)= \mathcal I'(u,v) (0,v)=0 \}.
\end{equation}
Precisely we characterize the seminodal least energy solution by the constrained
minimization problem
\[
\gamma_{\rm sn}=\inf \left\{ \mathcal I(u,v) : (u,v)\in \mathcal N_{\rm sn} \right\}.
\]
The converse case ($u$ sign-definite, $v$ sign-changing) can be handled with the
obvious adjustments.

Notice that $\mathcal N_{\rm nod}\subset \mathcal N_{\rm sn}$, therefore establishing
that $\gamma_{\rm sn}$ is attained  does not suffice  to ensure the existence of a
semi-nodal solution.

We mimic the reasoning of Section \ref{sec2} and introduce the auxiliary function
\[
\theta(a,b,c)=\mathcal I (a^{\frac{1}{q}} u^+
 - b^{\frac{1}{q}}u^-, c^{\frac{1}{q}} v),
 \]
and the sets
\begin{gather*}
 E_0 = \left\{ (u,v)\in E : M(u,v) \text{ is positive defined}\right\},
\\
  \mathcal O = \left\{ (u,v) \in E_0 :  M^{-1} W \in (0,\infty)^3\right\} ,
\end{gather*}
where now $M$ is the matrix
\[
M=\begin{pmatrix} \mu_1\|u^+\|_{2q}^{2q} & 0 & \beta \|u^+v\|_q^q \\
 0&  \mu_1\|u^-\|_{2q}^{2q} & \beta \|u^-v\|_q^q \\
 \beta \|u^+v\|_q^q & \beta \|u^-v\|_q^q &  \mu_2 \|v\|_{2q}^{2q}
\end{pmatrix} .
\]

The arguments of Lemmas \ref{lem:M>0}--\ref{prop:theta_q2}
give the following result.

\begin{lemma}\label{lem:snod}
For $|\beta|<\sqrt{\mu_1\mu_2/2}$,  $E_0=\{ (u,v) \in E : u^{\pm}, v \neq 0\}$.
For $\beta \le 0$,  $E_0$ contains any $(u,v)\in E$ such that
\begin{equation*}
\mu_1 \|u^{\pm}\|_{2q}^{2q} + \beta \|u^{\pm} v\|_{q}^q>0, \quad
\mu_2 \|v\|_{2q}^{2q} +\beta \|u v\|_{q}^q >0,
\end{equation*}
and in particular $\mathcal N_{\rm sn}\subset E_0$.

For $q=2$ and $0<\beta < \sqrt{\mu_1\mu_2/2}$, $\mathcal O$ is nonempty and
 $\mathcal N_{\rm sn}\subset \mathcal O$.

For $q\in (2,q_N)$ and $\beta<\sqrt{\mu_1\mu_2/2}$, or $q=2$ and $\beta \le 0$,
if $(u,v)\in  E_0$ then the function
$\theta$ has a unique maximum point in $ (0,\infty)^3$ characterized by the condition
$(a^{\frac{1}q} u^+ -b^{\frac{1}q} u^-, c^{\frac{1}q} v) \in \mathcal N_{\rm sn}$

If $q=2$ and $0<\beta< \sqrt{\mu_1\mu_2/2}$, the same holds true provided that
$(u,v)\in \mathcal O$.
\end{lemma}

Next, repeating the reasoning of Theorem \ref{thm:existence_ nodal} one shows
the following proposition.

\begin{proposition}\label{prop:seminodal1}
Let $q\in (2,q_N)$, for every $\beta < \sqrt{\mu_1\mu_2/2}$ the infimum $\gamma_{\rm sn}$
is attained by a function which solves \eqref{S}.
If $q=2$, there exists $\beta_0>0$ such that the same holds true for every
$\beta<\beta_0$.
\end{proposition}

The solution we  just constructed is not necessarily semi-nodal, because also $v$
could change sign. Next statement ensures that it is not the case.

\begin{proposition}\label{prop:seminodal2}
If $(u,v)\in \mathcal N_{\rm sn}$ realizes the minimum $c_{\rm sn}$, then it has
Morse index 3, $u$ has exactly two nodal zones and  $v$ has fixed sign.
\end{proposition}

\begin{proof}
Letting  $H:=E\cap \left(H^2(\Omega)\right)^2$ and using the function
$G:E_{\rm sn}\to\mathbb{R}^3$,
\[
G(u,v)=\begin{pmatrix}
\mathcal I'(u,v) \, (u^+,0) \\ \mathcal I'(u,v) \, (u^-,0) \\
\mathcal I'(u,v) \, (0,v)
 \end{pmatrix},
\]
one sees that $N_{\rm sn}\cap H$ is a $C^1$ manifold of codimension 3 in $H$.
So the same arguments of Proposition \ref{prop:Morseindex} yield that $(u,v)$
has Morse index 3.
Furthermore $u^{\pm} , v\neq 0$ by assumption. If $v$ changes sign or $u$ has more
than two nodal zones, it follows that the Morse index is at least 4, and it ends
the proof.
\end{proof}

Because $\mathcal N_{\rm nod} \subset \mathcal N_{\rm sn} \subset \mathcal N\subset
\mathcal N_{\rm full}$, a noteworthy consequence of this line of reasoning  is that
 all these inclusions are strict and we have the statement.

\begin{theorem}
$c_{\rm full} < c <c_{\rm sn} < c_{\rm nod}$.
\end{theorem}

\section{Symmetry breaking and more solutions with pre-assigned symmetries}\label{sec3}

Henceforth we take that $\Omega=B_R(0)$ is a ball.
The constrained minimization considered in the previous sections may be restricted
to the subspace consisting of radially symmetric functions
\[
E_{\rm rad}:=\{(u,v)\in E  : u(x)=u(|x|), \,
v(x)=v(|x|) \text{ for all } x\in \Omega \} ,
\]
and produces a pure vector, a nodal and a seminodal radial solutions, which have
respectively \emph{radial Morse index} equal to 2, 3 and 4.
By radial Morse index of a solution $(u,v)$ we mean the maximal dimension of a subspace
of $E_{\rm rad}$ where the quadratic form related to $\mathcal I''(u,v)$ is negative.
Observe that the radial Morse index can be less than the Morse index.

It remains to determine whether this provides genuinely new solutions or rather a
symmetry property of the least energy solutions.
One can find an answer, and observe symmetry breaking, by computing the exact Morse
index of radial solutions.
The issue is by itself interesting and deserves a detailed study. Here we argue by
perturbation and present an estimate in the case $\beta$ close to zero.

\begin{proposition}\label{prop:symmetry_breaking}
Let $N=2, 3$ and $2\le q<q_N$. There exist $b_1  , b_2> 0$ such that
\begin{enumerate}
 \item if $|\beta|<b_1$, a radial sign-changing least energy solution has
 Morse index greater or equal to $2N+2$.

\item if $|\beta|<b_2$, a radial seminodal least energy solution has Morse index
greater or equal to $N+ 2$.
\end{enumerate}
\end{proposition}

Comparing these estimates with Propositions \ref{prop:Morseindex} and
\ref{prop:seminodal2} proves the symmetry breaking stated by
Theorem \ref{thm:symmetry_breaking} .

Before proving Proposition \ref{prop:symmetry_breaking}, we check the uniform
convergence of solutions when $\beta$ vanishes.
To this end, we take $\lambda_i$, $\mu_i$ and $q$ as fixed, and  emphasize
the dependence by $\beta$ by writing $\mathcal I_{\beta}$ for the functional introduced
in \eqref{energy}  and $(u_{\beta}, v_{\beta})$ for a solution of \eqref{S}.
Recall that when $\beta=0$ system \eqref{S} decouples and gives rise to two
Lane-Emden problems
\begin{equation}  \label{eq:u}
\begin{gathered}
  -\Delta u +\lambda_1 u = \mu_1|u|^{2q-2} u \quad \text{in } \Omega , \\
  u=0 \quad \text{on } \partial\Omega,
 \end{gathered}
\end{equation}
and
\begin{equation}  \label{eq:v}
\begin{gathered}
 -\Delta v +\lambda_2 v = \mu_2|v|^{2q-2} v \quad \text{in } \Omega , \\
 v=0 \quad \text{on } \partial\Omega .
\end{gathered}
\end{equation}

\begin{lemma}\label{lem:stability}
For $\beta\neq 0$, let $(u_{\beta}, v_{\beta})$ a nodal (respectively, semi-nodal)
least energy radial solution to \eqref{S}.
There are $u_0$ and $v_0$, nontrivial  radial solutions to \eqref{eq:u} and
\eqref{eq:v}, such that  $u_{\beta}\to u_0$ and  $v_{\beta}\to v_0$  uniformly
as $\beta\to0$, up to an extracted sequence.
Furthermore if $(u_{\beta}, v_{\beta})$ are nodal (resp., semi-nodal), the same
holds for $u_0$ and $v_0$.
 \end{lemma}

\begin{proof}
We take that for every $\beta\neq 0$, $(u_\beta, v_\beta)$ is a radial sign-changing
least energy solution.
By Proposition \ref{prop:existence_radial} we know that
$(u_\beta, v_\beta)\in \mathcal N_{\rm nod} \cap E_{\rm rad}$ and
\[
  \mathcal I_{\beta}(u^{\rm nod}_{\beta}, v^{\rm nod}_{\beta})
  = \min \left\{ \mathcal I_{\beta}(u,v) : (u,v) \in \mathcal N_{\rm nod}
  \cap E_{\rm rad}\right\} .
\]
The  arguments used in the proof of Theorem \ref{thm:existence_ nodal},
to obtain \eqref{bound_nod}, proves that there exist $C$ not depending by $\beta$
such that
\begin{equation} \label{bound_unif}
\|u_{\beta}\|_{2q},  \|v_{\beta}\|_{2q} \le C,
\end{equation}
and since $(u_{\beta},v_{\beta})\in \mathcal N$ it follows readily that
\begin{equation} \label{bound_unif_E}
 \|u_{\beta}\|_{E_1},  \|v_{\beta}\|_{E_2} \le C
\end{equation}
for $|\beta|<b$.
Therefore by standard arguments $u_{\beta}$ and $v_{\beta}$ converge weakly in $H^1_0$,
strongly in $L^{2q}(\Omega)$ and pointwise a.e.~to radial solutions $u_0$ and $v_0$
of \eqref{eq:u}, \eqref{eq:v}.
Reasoning as in the step 2 of the proof of Theorem \ref{thm:existence_ nodal} one
deduces from \eqref{bound_unif} that $u_0$ $v_0$ are sign-changing
(in particular, non-trivial).

If $(u_\beta, v_\beta)$ are semi-nodal, let us say that $u_{\beta}$ are sign-changing
and $v_{\beta}>0$ on $\Omega$, the same reasoning yields that $u_0$ changes sign and
$v_0$ is not identically zero.
Eventually $v_0\ge0$ by the pointwise convergence, and  Hopf's boundary Lemma assures
that $v_0>0$ on $\Omega$.

It remains to see that the convergence is uniform, indeed.
First we point out that radial weak solutions to \eqref{S} are indeed classical
and integrating the equation in \eqref{S} gives
\begin{gather}\label{u'}
 u_{\beta}'(r)=-\frac{1}{r^{N-1}} \int_0^r \rho^{N-1}
  \left( \mu_1|u_{\beta}|^{2q-2}+\beta |u_{\beta}|^{q-2}|v_{\beta}|^q -\lambda_1
  \right) u_{\beta} d\rho , \\
\label{v'}
  v_{\beta}'(r)=-\frac{1}{r^{N-1}} \int_0^r \rho^{N-1}
  \left( \mu_2|v_{\beta}|^{2q-2}+\beta |u_{\beta}|^{q}|v_{\beta}|^{q-2} -\lambda_2 \right)
   v_{\beta} d\rho, \\
 \label{uv}
  u_{\beta}(r)= - \int_r^R  u_{\beta}'(\rho) d\rho , \quad  v_{\beta}(r)
  = - \int_r^R  v_{\beta}'(\rho) d\rho .
 \end{gather}

Estimating \eqref{u'} and \eqref{v'} by Holder's inequality and \eqref{bound_unif} gives
\begin{equation} \label{u'v'Holder}
|u_{\beta}'(r)|, | v_{\beta}'(r)|
\le C r^{1-N} \Big(r^{\frac{N}{2q}} + r^{N\frac{2q-1}{2q}} \Big) .
\end{equation}
If $N=2$,  inserting \eqref{u'v'Holder} inside \eqref{uv} yields
$|u_{\beta}(r)|, | v_{\beta}(r)|\le C$, which in turn, together with \eqref{u'}
and \eqref{v'},
ensures that $|u'_{\beta}(r)|, | v'_{\beta}(r)|\le C$ and concludes the proof.

If $N= 3$, instead, it is needed to start a bootstrap argument by the so called Radial
Lemma due to Strauss \cite{strauss}: if $w$ is any radial function in $H^1_0(\Omega)$,
then
\begin{equation} \label{rl}
|w(r)|\le \frac{\|\nabla w\|_{2} }{\sqrt{N-2}} \, r^{-\frac{N-2}2}.
\end{equation}
To simplify notations, we write $p=2q-1$.
Estimating \eqref{u'} and \eqref{v'} by \eqref{rl} and \eqref{bound_unif_E} gives
\begin{equation} \label{u'v'rl}
 |u_{\beta}'(r)|, | v_{\beta}'(r)|\le C \Big( 1+ r^{1-p\frac{N-2}{2}} \Big) .
\end{equation}
 If $1-p\frac{N-2}{2} >-1$, the proof ends as for  $N=2$.
Also when $1-p\frac{N-2}{2} =-1$, that is $p= \frac{4}{N-2}$, by
 plugging \eqref{u'v'rl} into \eqref{uv} we obtain
\[
|u_{\beta}(r)|, | v_{\beta}(r)|\le C \left( 1+ \log r\right) ,
\]
and going back to \eqref{u'} and \eqref{v'}, we ave
\[
|u_{\beta}'(r)|, | v_{\beta}'(r)|
\le C  \Big(1+ r^{1-N}\int_0^r \rho^{N-1} |\log\rho|^{\frac4{N-2}}d\rho \Big) \le C ,
\]
which ensures uniform convergence.
Otherwise if $1-p\frac{N-2}{2} <-1$, we put \eqref{u'v'rl} into \eqref{uv} and obtain
\begin{equation} \label{uvrl}
|u_{\beta}(r)|, | v_{\beta}(r)|\le C \left( 1+ r^{2-p\frac{N-2}{2}} \right) ,
\end{equation}
which in turn implies
\[
|u_{\beta}'(r)|, | v_{\beta}'(r)|\le C \left( 1+ r^{1+2p -p^2\frac{N-2}{2}} \right) .
\]
If $1+2p -p^2\frac{N-2}{2} \ge -1$ the proof ends like in the previous step,
 otherwise we iterate the argument  and end up with
\begin{align*}%\label{u'v'rl3}
|u_{\beta}'(r)|, | v_{\beta}'(r)|\le C \left( 1+ r^{1+2p +2p^2-p^3\frac{N-2}{2}} \right) . \end{align*}
In that way, after $n$ steps we get
\[
|u_{\beta}'(r)|, | v_{\beta}'(r)|\le C \left( 1+ r^{-1+\gamma_n} \right) ,
\]
for
\[
\gamma_n= 2\sum_{k=0}^n p^k-p^{n+1}\frac{N-2}{2}
= \Big(\frac{2}{p-1} - \frac{N-2}{2} \Big) p^{n+1}-  \frac{2}{p-1} ,
\]
and the proof is concluded if $\gamma_n\ge 0$  after a finite number of steps.
But this must hold because $\frac{2}{p-1} = \frac{1}{q-1} >  \frac{N-2}{2}$.
\end{proof}

\begin{proof}[Proof of Proposition \ref{prop:symmetry_breaking}]
We prove in full detail part (1).
Using the same notations of Lemma \ref{lem:stability}, it is know by
\cite[Theorem 1.1]{AftPac} that  $u_0$ (resp., $v_0$) has Morse index greater or
equal to $N+1$, as a solution to \eqref{eq:u} (resp., \eqref{eq:v}).
This means that there are $N+1$ linearly independent functions
$\phi_i\in H^1_0(\Omega)$  that solve
\begin{equation}\label{autofunzu}
  -\Delta \phi_i +\lambda_1\phi_i= (2q-1)\mu_1|u_0|^{2q-2}\phi_i + \nu_i \phi_i
\end{equation}
in $\Omega$, with $\nu_1\le \nu_2\dots \le\nu_{N+1}<0$ and
 $N+1$ linearly independent functions $\psi_i\in H^1_0(\Omega)$ that solve
\begin{equation}\label{autofunzv}
  -\Delta \psi_i +\lambda_2\psi_i= (2q-1)\mu_2|v_0|^{2q-2}\psi_i + \sigma_i \psi_i
\end{equation}
in $\Omega$, with $\sigma_1\le \sigma_2\dots \le\sigma_{N+1}<0$.

Let $W_1 , W_2$ be the $(N\!+\!1)$--dimensional subspaces of $H^1_0(\Omega)$ spanned  by
$\phi_1,\dots \phi_{N+1}$ and $\psi_1, \dots \psi_{N+1}$, respectively.

The claim follows by checking that, if $|\beta|$ is small, then the quadratic form
related to $\mathcal I_{\beta}''(u_{\beta}, v_{\beta})$ is negative on $W_1\times W_2$.
Recall that by \eqref{autofunzu} and \eqref{autofunzv} we have
\begin{gather*}
\|\phi\|_{E_1} \le (2q-1)\mu_1 \int_{\Omega}|u_0|^{2q-2} \phi^2 +\nu_{N+1} \|\phi\|_2^2
\quad \text{for every } \phi\in W_1,
  \\
\|\psi\|_{E_2} \le (2q-1)\mu_2 \int_{\Omega}|v_0|^{2q-2} \psi^2 +\sigma_{N+1}
 \| \psi\|_2^2 \quad \text{for every } \psi\in W_2 .
 \end{gather*}
Hence
\begin{align*}
&\langle \mathcal I_{\beta}''(u_{\beta}, v_{\beta}) (\phi,\psi), (\phi,\psi)\rangle \\
&=\|\phi\|_{E_1}^2+\|\psi\|_{E_1}^2
 - (2q-1)\mu_1 \int_{\Omega}|u_{\beta}|^{2q-2} \phi^2  - (2q-1)\mu_2
 \int_{\Omega}|v_{\beta}|^{2q-2} \psi^2
 \\
&\quad- (q-1)\beta\int_{\Omega}\left(|u_{\beta}|^{q-2}|v_{\beta}|^q \phi^2 + |u_{\beta}|^{q}|v_{\beta}|^{q-2} \psi^2 \right)
 - 2q\beta \int_{\Omega}|u_{\beta}v_{\beta}|^{q-2} u_{\beta}v_{\beta} \phi\psi
 \\
&\le  \nu_{N+1} \|\phi\|_2^2+\sigma_{N+1} \| \psi\|_2^2
 + (2q-1)\mu_1 \int_{\Omega}\left(|u_0|^{2q-2} - |u_{\beta}|^{2q-2} \right)\phi^2 \\
&\quad +  (2q-1)\mu_2 \int_{\Omega}\left(|v_0|^{2q-2} - |v_{\beta}|^{2q-2} \right)
 \psi^2  - (q-1)\beta\int_{\Omega}\left(|u_{\beta}|^{q-2}|v_{\beta}|^q \phi^2
  + |u_{\beta}|^{q}|v_{\beta}|^{q-2} \psi^2 \right) \\
&\quad - 2q\beta \int_{\Omega}|u_{\beta}v_{\beta}|^{q-2} u_{\beta}v_{\beta} \phi\psi
 \intertext{and thanks to the uniform convergence in Lemma \ref{lem:stability}}
&\le \left(\max\{\mu_{N+1}, \sigma_{N+1}\} + o(1)\right)
\left( \|\phi\|_2^2 + \|\psi\|_2^2\right)
\end{align*}
where $o(1)\to 0 $ as $\beta\to 0$.
\end{proof}

\subsection{Other type of symmetries and multiplicity of solutions}
Here $\Omega\subset \mathbb{R}^2$ is a disc.
We write $r$ and $\theta$ for the usual polar coordinates, and  for every
$k\in \mathbb Z_+$ we introduce the set of $k$-symmetric functions
\[
H_{k}:=\{(u,v)\in \left(H^1_0(\Omega)\right)^2  : \text{ $u$ and $v$ are
even and $\frac{2\pi}{k}$-periodic w.r.t. $\theta$} \} .
\]
For a \emph{$k$-symmetric} solution, we denote by $\mathrm m_k$ its $k$ Morse index,
i.e.\ the maximal dimension of a subspace of $H_k$ where the second derivative of
$\mathcal I$ gives rise to a negative defined quadratic form.
By constrained minimization in the symmetric Nehari set
$\mathcal N_{\rm nod}\cap H_k$ it is straightforward to check the  following.

 \begin{proposition}\label{prop:existence_radial}
Let $\Omega\subset \mathbb{R}^2$ be a disc, $q\ge 2$, then there exist $\bar \beta_k>0$,
such that for $\beta<\beta_k$ the problem \eqref{S} has
a (component-wise) sign-changing  $k$-symmetric  solution which has the least
energy among all sign-changing and $k$-symmetric solutions and  has $k$ Morse index 4.
 \end{proposition}

Relying on the convergence in Lemma \ref{lem:stability} and the results about
Lane-Emden equation in \cite{GlaIan23Nodea}
we can estimate the $k$-Morse index when $\beta$ is close to zero, thus proving
 Theorem \ref{thm:ksym}.

\begin{proof}[Proof of Theorem \ref{thm:ksym}]
If a sign-changing $k$-symmetric least energy solution to \eqref{S}
is radial, then it has to be a sign-changing radial ground state, because the set
of component-wise radial functions is contained in  $H_{k}$.
We prove this cannot happen by showing that the $k$ Morse index of a sign-changing radial
ground state is at least $6$, provided that $q$ is large  and $\beta$ is close to $0$.
To this aim, we denote by $(u_{\beta}, v_{\beta})$ a sign-changing radial solution.
For $\beta\to 0$,  Lemma \ref{lem:stability} states that  $u_{\beta}$ and $v_{\beta}$
converge uniformly to $u_0$ and $v_0$,  which are sign-changing radial  solutions
to \eqref{eq:u} and \eqref{eq:v}, respectively.
It is known by \cite[Proposition 6.5]{GlaIan23Nodea} that  the $k$ Morse index of both
$u_0$ and $v_0$ (as solutions of \eqref{eq:u} and \eqref{eq:v}, respectively) is not
less than $3$ when $q$ is large.
Hence reasoning as in the proof of Proposition \ref{prop:symmetry_breaking} one sees
that the quadratic form related to $ \mathcal I_{\beta}''(u_{\beta}, v_{\beta})$
is negative on a subset of $H_k$ of dimension $6$, if $\beta$ is close to $0$.
\end{proof}

\subsection*{Acknowledgments}
This work was  supported by Gruppo Nazionale per l'Analisi Matematica,
la Probabilit\`a e le loro Applicazioni (GNAMPA) of the Istituto Nazionale di
Alta Matematica (INdAM) and by the
Geometric-Analytic Methods for PDEs and Applications (GNAMPA)
project CUP I53D23002420006- funded by the European Union -
Next Generation EU within the PRIN 2022 program
(D.D. 104 - 02/02/2022 Ministero dell'Universit\`a e della Ricerca).
This manuscript reflects only the author's views and opinions and the Ministry
cannot be considered responsible for them.

 \begin{thebibliography}{10}

\bibitem{AftPac}
 Aftalion, A.; Pacella, F.;
 Qualitative properties of nodal solutions of semilinear
 elliptic equations in radially symmetric domains.
 \emph{C. R. Acad. Sci.}, 339 (2004), 339-344.

\bibitem{AmaGla20NARWA}
Amadori, A.L.; Gladiali, F.;
 On a singular eigenvalue problem and its applications in computing the Morse index of solutions to semilinear PDE’s,
 \emph{Nonlinear Analysis: Real World Applications},
 Volume 55,  103133, (2020),
 DOI: 10.1016/j.nonrwa.2020.103133.

\bibitem{Ama24JEPE}
Amadori, A.L.;
 Some remarks about the Morse index for convex Hamiltonian systems
\emph{Journal of Elliptic and Parabolic Equations}, 10  (2024) (2), pp. 1255 - 1274,
 DOI: 10.1007/s41808-024-00295-3

 \bibitem{BW_JDE06}
 Bartsch, T.; Wang, Z. Q.;
 Note on Ground States of Nonlinear Schr\"{o}dinger Systems
\emph{Journal of Partial Differential Equations}, Vol. 19 (2006), Iss. 3 : pp. 200-207.

\bibitem{CLZ_JDE13}
Chen, Z.; Lin, C.-S.; Zou, W.;
Multiple sign-changing and semi-nodal solutions for coupled Schr\"{o}dinger equations,
\emph{Journal of Differential Equations}(2013)
DOI: 10.1016/j.jde.2013.08.009.

\bibitem{CLZ_ASNSP16}
Chen, Z.; Lin, C.-S.; Zou, W.;
Infinitely many sign-changing and semi-nodal solutions for a nonlinear Schr\"{o}dinger
system, \emph{Annali della Scuola normale superiore di Pisa - Classe di scienze} (2016)
DOI: 10.48550/arXiv.1212.3773

\bibitem{CS_CCM23}
Clapp, M.; Soares, M.;
Coupled and uncoupled sign-changing spikes of singularly perturbed elliptic systems,
 \emph{Communications in Contemporary Mathematics} (2023) DOI: 10.1142/S0219199722500481

\bibitem{CTV_AIHP}
Conti, M.; Terracini, S.;  Verzini, G.;
 Nehari's problem and competing species systems.
\emph{Ann. Inst. H. Poincar\'e Anal. Non Lin\'eaire} (2002),
DOI: 10.1016/s0294-1449(02)00104-x

\bibitem{BoseEinstein}
Esry, B. D.; Greene, H.H.; Burke Jr. J. P.; Bohn, J. L.;
Hartree-Fock theory for double condensates, \emph{Phys. Rev. Lett.} 78 (1997) 3594-3597.

\bibitem{DamPac}
 Damascelli, L.; Pacella, F.;
Symmetry Results for Cooperative Elliptic Systems via Linearization
\emph{SIAM Journal on Mathematical Analysis} (2013) 45:3, 1003-1026,
DOI: 10.1137/110853534

\bibitem{GNN}
Gidas, B.; Ni, W. M.; Nirenberg, L.;
Symmetry and related properties via the maximum
principle. \emph{Comm. Math. Phys.}, 68 (1979) (3), 209--243.

\bibitem{GlaIan23Nodea}
Gladiali, F.; Ianni, I.;
Quasi-radial solutions for the Lane-Emden problem in the ball
\emph{Nonlinear Differential Equations and Applications}, 27 (2), art. no. 13 (2020)
DOI: 10.1007/s00030-020-0616-0

\bibitem{LLW_CV15}
Liu, J.; Liu, X.; Wang, Z.;
Multiple mixed states of nodal solutions for nonlinear Schr\"{o}dinger systems.
\emph{Calc. Var.}, (2015), DOI: 10.1007/s00526-014-0724-y

\bibitem{LW_AIHP05}
Lin, T.; Wei, J.;
Spikes in two coupled nonlinear Schr\"{o}dinger equations.
\emph{Ann. Inst. H.   Poincar\'e Anal. Non Lin\'eaire 22} (2005),
  DOI: 10.1016/j.anihpc.2004.03.004

\bibitem{MMP_JDE06}
Maia, L.; Montefusco, E.; Pellacci, B.;
Positive solutions for a Weakly Coupled nonlinear Schr\"odinger system,
\emph{Journal of Differential Equations} (2006),
DOI: 10.1016/j.jde.2006.07.002.

\bibitem{MMP_CCM08}
Maia, L.; Montefusco, E.; Pellacci, B.;
Infinitely many nodal solutions for a Weakly Coupled nonlinear Schr\"odinger system,
\emph{Commun.Contemp.Math.}, (2008),  DOI: 10.1142/S0219199708002934

\bibitem{Miranda}
 Miranda, C.;
Un'osservazione su un teorema di Brouwer.
\emph{Boll. Un. Mat. Ital.} 2(1940),(3), 5--7,

\bibitem{PacWet}
Pacella, F.;  Weth, T.;
Symmetry Results for Solutions of Semilinear Elliptic
Equations via Morse index, \emph{Proc. AMS}, 135 (2007) (6), pp 1753-1762.

\bibitem{SW_DCDS15}
Sato, Yohei K.;   Wang, Z.-Q.;
On the least energy sign-changing solutions for a nonlinear elliptic system.
\emph{Discrete and Continuous Dynamical Systems} (2015)
DOI: 10.3934/dcds.2015.35.2151

\bibitem{Song_JDE25}
Song, Y.;
Infinitely many sign-changing and semi-nodal solutions to Schr\"{o}dinger systems
with mixed competition and cooperation terms,
\emph{Journal of Differential Equations} (2025), DOI: 10.1016/j.jde.2024.12.011

\bibitem{strauss}
Strauss, W. A.;
Existence of solitary waves in higher dimensions.
\emph{Commun. Math. Phys.}, 55, 149-162 (1977), DOI: 10.1007/BF01626517

\bibitem{TT_AIHP12}
Tavares, H.; Terracini, S.;
Sign-changing solutions of competition-diffusion elliptic systems and optimal
partition problems,
\emph{Annales de l'Institut Henri Poincare C, Analyse non linaire} (2012),
 DOI: 10.1016/j.anihpc.2011.10.006.

\bibitem{Z_BVP24}
Zhang, J.;
Sign-changing solutions for coupled Schr\"odinger system,
\emph{Bound Value Probl} (2024), DOI: 10.1186/s13661-024-01881-z

\end{thebibliography}

\end{document}

